\chapter{Variations of Energy}
\label{chap:dc-variations-of-energy}

Faithful transcription of do Carmo, \emph{Riemannian Geometry}. Each numbered
statement keeps do Carmo's number, kind and (name); proofs keep his explicit
cross-references ("by \cref{lem:dc-ch9-2-3}", "Theor. 3.7, Chap. 3") so the dependency graph
is recoverable from the text.

\section{Introduction}

In Chapter 3, geodesics were defined as curves with zero acceleration, and were
characterized by the fact that they locally minimize arc length. This chapter
presents a further characterization of a geodesic as a "solution of a variational
problem", adapting to Differential Geometry concepts and techniques from the
Calculus of Variations. The fundamental point of the chapter is the calculation
of the formula for the second variation of the energy of a geodesic (Section 2).
In Section 3 two geometric applications are made: the Theorem of Bonnet–Myers,
stating that a complete manifold whose curvature is positive and does not approach
zero is compact, with diameter estimated in terms of bounds on the curvature; and
an extension, due to A. Weinstein, of a theorem of Synge asserting the simple
connectivity of a compact, orientable manifold of even dimension whose curvature
is positive. Together with the theorem of Hadamard, these applications investigate
the influence of curvature on the topology of Riemannian manifolds, culminating in
the Sphere Theorem (Chap. 13).

\section{Formulas for the first and second variations of energy}

We make precise the idea of "neighboring curves" to a given curve.

\begin{definition}[variation]
  \label{def:dc-ch9-2-1}
  \dcref{ch9:2.1}
  Let $c:[0,a]\to M$ be a piecewise differentiable curve in a manifold $M$. A
  \emph{variation} of $c$ is a continuous mapping $f:(-\varepsilon,\varepsilon)\times[0,a]\to M$
  such that:
  - a) $f(0,t)=c(t)$, $t\in[0,a]$,
  - b) there exists a subdivision of $[0,a]$ by points
  $0=t_0<t_1<\cdots<t_{k+1}=a$, such that the restriction of $f$ to each
  $(-\varepsilon,\varepsilon)\times[t_i,t_{i+1}]$, $i=0,1,\dots,k$, is differentiable.

  A variation is said to be \emph{proper} if $f(s,0)=c(0)$ and $f(s,a)=c(a)$, for all
  $s\in(-\varepsilon,\varepsilon)$. If $f$ is differentiable, the variation is said to be
  \emph{differentiable}. For each $s\in(-\varepsilon,\varepsilon)$, the parametrized curve
  $f_s:[0,a]\to M$ given by $f_s(t)=f(s,t)$ is called a \emph{curve in the variation}, so
  a variation determines a family $f_s(t)$ of neighboring curves of $f_0(t)=c(t)$.
  The variation is proper exactly when the curves in this family have the same
  initial point $c(0)$ and endpoint $c(a)$.
  The parametrized differentiable curve given by $f_t(s)=f(s,t)$, $t$ fixed, is a
  \emph{transversal curve of the variation}. The velocity vector of a transversal curve
  at $s=0$, defined by $V(t)=\frac{\partial f}{\partial s}(0,t)$, is a (piecewise
  differentiable) vector field along $c(t)$ and is called the \emph{variational field} of
  $f$.
\end{definition}

\begin{proposition}[existence of a variation with given variational field]
  \label{prop:dc-ch9-2-2}
  \dcref{ch9:2.2}
  Given a piecewise differentiable field $V(t)$, along the piecewise differentiable
  curve $c:[0,a]\to M$, there exists a variation $f:(-\varepsilon,\varepsilon)\times[0,a]\to M$
  of $c$, such that $V(t)$ is the variational field of $f$; in addition, if
  $V(0)=V(a)=0$, it is possible to choose $f$ as a proper variation.
\end{proposition}
\begin{proof}
  Since $c([0,a])\subset M$ is compact it is possible to find a $\delta>0$
  such that $\exp_{c(t)}$ is well-defined for all $v\in T_{c(t)}M$ with $|v|<\delta$:
  for each $c(t)$ consider a totally normal neighborhood $W_t$ of $c(t)$ and the
  number $\delta_t>0$ associated to it (Theor. 3.7, Chap. 3); a finite subcover
  $W_1,\dots,W_n$ of $c([0,a])$ gives $\delta=\min(\delta_1,\dots,\delta_n)$.
  Consider $N=\max_{t\in[0,a]}|V(t)|$, $\varepsilon<\frac{\delta}{N}$ and define
  $f(s,t)=\exp_{c(t)}sV(t)$. Since $\exp_{c(t)}sV(t)=\gamma(1,c(t),sV(t))$ and the
  geodesic depends differentiably on the initial conditions, $f$ is piecewise
  differentiable, and $f(0,t)=c(t)$. The variational field of $f$ is

  $$
  \frac{\partial f}{\partial s}(0,t)=\frac{d}{ds}(\exp_{c(t)}sV(t))\Big|_{s=0}=(d\exp_{c(t)})_0 V(t)=V(t),
  $$

  and if $V(0)=V(a)=0$ then $f$ is proper. $\square$

  To compare the arc length of $c$ with that of neighboring curves we define
  $L:(-\varepsilon,\varepsilon)\to\mathbf{R}$ by

  $$
  L(s)=\int_0^a\Big|\frac{\partial f}{\partial t}(s,t)\Big|\,dt,
  $$

  the length of the curve $f_s(t)$. It is more convenient to work with the *energy
  function* $E(s)$ given by

  $$
  E(s)=\int_0^a\Big|\frac{\partial f}{\partial t}(s,t)\Big|^2\,dt,\qquad s\in(-\varepsilon,\varepsilon).
  $$

  For a curve $c:[0,a]\to M$ set

  $$
  L(c)=\int_0^a\Big|\frac{dc}{dt}\Big|\,dt,\qquad E(c)=\int_0^a\Big|\frac{dc}{dt}\Big|^2\,dt.
  $$

  Putting $f\equiv 1$ and $g=|\frac{dc}{dt}|$ in the Schwarz inequality
  $(\int_0^a fg\,dt)^2\le\int_0^a f^2\,dt\cdot\int_0^a g^2\,dt$, we obtain

  $$
  L(c)^2\le a\,E(c),
  $$

  with equality if and only if $g$ is constant, that is, if and only if $t$ is
  proportional to arc length.
\end{proof}

The energy functional and the Schwarz comparison introduced in the proof above are used
on their own throughout the chapter (in \cref{lem:dc-ch9-2-3} and again at the end of the
proof of \cref{thm:dc-ch9-3-7}), so they are recorded here as separate nodes.

\begin{definition}[energy of a curve]
  \label{def:dc-ch9-2-2-energy}
  \dcref{ch9:2.2}
  \lean{Riemannian.DCEnergy, Riemannian.dcSpeed, Riemannian.dcArcLength_eq_integral_dcSpeed,
    Riemannian.dcEnergy_eq_integral_dcSpeed_sq}
  \leanok
  \uses{def:dc-ch1-2-9}
  For a curve $c:[a,b]\to M$ the \emph{speed} is $|dc/dt|=\langle dc/dt,dc/dt\rangle^{1/2}$
  and the \emph{energy} of the segment $c|[a,b]$ is
  $$
  E(c)=\int_a^b\Big|\frac{dc}{dt}\Big|^2\,dt,
  $$
  the companion of the arc length $L(c)=\int_a^b|dc/dt|\,dt$ of Def. 2.9, Chap. 1. Both are
  integrals of the speed: $L(c)=\int_a^b|dc/dt|\,dt$ and $E(c)=\int_a^b|dc/dt|^2\,dt$.
\end{definition}

\begin{lemma}[Schwarz inequality $L(c)^2\le a\,E(c)$]
  \label{lem:dc-ch9-2-2-schwarz}
  \dcref{ch9:2.2}
  \lean{Riemannian.dcArcLength_sq_le_mul_dcEnergy, Riemannian.dcArcLength_sq_eq_iff,
    Riemannian.sq_integral_le_sub_mul_integral_sq, Riemannian.sq_integral_eq_iff_ae_eq_const}
  \leanok
  \uses{def:dc-ch9-2-2-energy}
  For a curve $c:[a,b]\to M$ whose speed and squared speed are integrable,
  $$
  L(c)^2\le (b-a)\,E(c),
  $$
  with equality if and only if the speed $|dc/dt|$ is (almost everywhere) constant, that
  is, if and only if $t$ is proportional to arc length.

  (Only integrability is assumed, not continuity of the speed: the curves of
  \cref{def:dc-ch9-2-1} are merely \emph{piecewise} differentiable, so $|dc/dt|$ may jump
  at the subdivision points $t_i$, and a continuous-speed hypothesis would exclude
  precisely the curves this chapter is about.)
\end{lemma}
\begin{proof}
  \leanok
  \uses{def:dc-ch9-2-2-energy}
  Let $f=|dc/dt|$ and let $m=\frac{1}{b-a}\int_a^b f$ be its mean. Expanding the square,
  $$
  0\le\int_a^b (f-m)^2\,dt=\int_a^b f^2\,dt-2m\int_a^b f\,dt+m^2(b-a)=E(c)-\frac{L(c)^2}{b-a},
  $$
  which is the inequality. Equality holds exactly when $\int_a^b(f-m)^2\,dt=0$, and a
  non-negative integrand has vanishing integral precisely when it vanishes almost
  everywhere; that is, exactly when $f=m$ almost everywhere. $\square$
\end{proof}

\begin{lemma}[minimizing geodesics minimize energy]
  \label{lem:dc-ch9-2-3}
  \dcref{ch9:2.3}
  \uses{cor:dc-ch3-3-9, lem:dc-ch9-2-2-schwarz, lem:dc-ch9-2-3-inequality,
    lem:dc-ch9-2-3-equality-case, lem:dc-ch9-2-3-arclength-bridge}
  Let $p,q\in M$ and let $\gamma:[0,a]\to M$ be a minimizing geodesic joining $p$ to
  $q$. Then, for all curves $c:[0,a]\to M$ joining $p$ to $q$,

  $$
  E(\gamma)\le E(c)
  $$

  with equality holding if and only if $c$ is a minimizing geodesic.
\end{lemma}
\begin{proof}
  From the considerations above,

  $$
  aE(\gamma)=(L(\gamma))^2\le(L(c))^2\le aE(c),
  $$

  which proves the first assertion. If equality holds, then $(L(c))^2=aE(c)$, so the
  parameter of $c$ is proportional to arc length, and $L(\gamma)=L(c)$, so $c$ is a
  minimizing geodesic (see \cref{cor:dc-ch3-3-9}). The converse is obvious.
  $\square$
\end{proof}

The steps of that chain are recorded below as separate nodes, for two reasons. First, the
chain is assembled from pieces that are each reused (the geodesic equality case is used
again in \cref{thm:dc-ch9-3-1}, and the $L$--$E$ comparison at the end of
\cref{thm:dc-ch9-3-7}). Second --- and this is why \cref{lem:dc-ch9-2-3} itself is
\emph{not} marked \leanok\ --- do Carmo's statement quantifies over "all curves $c$ joining
$p$ to $q$" and closes with "$c$ is a minimizing geodesic", both of which presuppose that
minimality can be \emph{compared} against the arc length $L$ of Def. 2.9, Chap. 1
(\cref{def:dc-ch1-2-9}). The first node below supplies exactly that comparison; the
remaining ones are the chain proper. Only the final appeal to \cref{cor:dc-ch3-3-9} --- "so
$c$ is a geodesic" --- is still open, on two counts: that corollary assumes arc length
proportional to the parameter \emph{pointwise}, whereas the equality case yields it only
almost everywhere; and its minimality hypothesis is quantified over \emph{every} competing
curve, whereas the equality case gives $L(\gamma)=L(c)$ for the single $c$ at hand.

\begin{lemma}[arc length as a path length]
  \label{lem:dc-ch9-2-3-arclength-bridge}
  \dcref{ch9:2.3}
  \lean{Riemannian.ofReal_dcArcLength_eq_pathELength,
    Riemannian.enorm_mfderiv_eq_ofReal_dcSpeed,
    Riemannian.dcArcLength_eq_toReal_pathELength,
    Riemannian.dcArcLength_le_of_pathELength_le}
  \leanok
  \uses{def:dc-ch1-2-9, def:dc-ch9-2-2-energy}
  For a curve $c:[a,b]\to M$ whose speed is integrable, the arc length $L(c)$ of Def. 2.9,
  Chap. 1 --- an integral of $|dc/dt|$ against $dt$ --- agrees with the length of $c$
  measured as a path in the metric space $M$. Consequently a comparison of path lengths
  between two curves with integrable speeds is a comparison of their arc lengths $L$.

  (The integrability hypothesis is not removable: a path length is allowed to be infinite,
  while $L(c)=\int_a^b|dc/dt|\,dt$ of a non-integrable speed is not, so for such a curve the
  two genuinely differ.)
\end{lemma}
\begin{proof}
  \leanok
  \uses{def:dc-ch1-2-9}
  The two integrands agree pointwise: the norm of the velocity is $|dc/dt|$ by definition of
  the metric on the tangent spaces. The two notions therefore differ only in the mode of
  integration --- an integral of a real-valued function against $dt$, versus a lower integral
  of the corresponding non-negative extended-real function --- and for a non-negative
  integrable function the two agree. $\square$
\end{proof}

\begin{lemma}[a geodesic attains equality in the Schwarz comparison]
  \label{lem:dc-ch9-2-3-geodesic-equality}
  \dcref{ch9:2.3}
  \lean{Riemannian.dcArcLength_sq_eq_mul_dcEnergy_of_isGeodesicOn,
    Riemannian.dcSpeed_eq_of_isGeodesicOn}
  \leanok
  \uses{lem:dc-ch9-2-2-schwarz, lem:dc-ch3-2-1-constspeed, def:dc-ch9-2-2-energy}
  For a geodesic $\gamma:[a,b]\to M$ with integrable speed and squared speed,
  $$
  L(\gamma)^2 = (b-a)\,E(\gamma).
  $$
  This is do Carmo's "$aE(\gamma)=(L(\gamma))^2$", the first step of the chain of
  \cref{lem:dc-ch9-2-3}.
\end{lemma}
\begin{proof}
  \leanok
  \uses{lem:dc-ch9-2-2-schwarz, lem:dc-ch3-2-1-constspeed}
  A geodesic has constant speed (\cref{lem:dc-ch3-2-1-constspeed}), so its parameter is
  proportional to arc length, which is precisely the equality case of
  \cref{lem:dc-ch9-2-2-schwarz}.

  (Constant speed is obtained by differentiating $\langle\gamma',\gamma'\rangle$ and is
  therefore available only on an \emph{open} interval of times; the equality case of
  \cref{lem:dc-ch9-2-2-schwarz} asks for constancy only almost everywhere, and the two
  endpoints form a null set, so nothing is lost.) $\square$
\end{proof}

\begin{lemma}[a minimizing geodesic minimizes energy: the inequality]
  \label{lem:dc-ch9-2-3-inequality}
  \dcref{ch9:2.3}
  \lean{Riemannian.dcEnergy_le_of_dcArcLength_le}
  \leanok
  \uses{lem:dc-ch9-2-3-geodesic-equality, lem:dc-ch9-2-2-schwarz, def:dc-ch9-2-2-energy}
  Let $\gamma:[a,b]\to M$ be a geodesic and $c:[a,b]\to M$ a curve with $L(\gamma)\le L(c)$
  (both with integrable speed and squared speed). Then $E(\gamma)\le E(c)$.

  (Minimality is assumed in the form $L(\gamma)\le L(c)$, which is the only consequence of
  "$\gamma$ is minimizing" that the proof uses; \cref{lem:dc-ch9-2-3-arclength-bridge} is
  what converts a metric minimality statement into it.)
\end{lemma}
\begin{proof}
  \leanok
  \uses{lem:dc-ch9-2-3-geodesic-equality, lem:dc-ch9-2-2-schwarz}
  do Carmo's chain verbatim:
  $(b-a)E(\gamma)=L(\gamma)^2\le L(c)^2\le (b-a)E(c)$, the first equality by
  \cref{lem:dc-ch9-2-3-geodesic-equality}, the middle step by hypothesis (both lengths being
  non-negative), and the last by \cref{lem:dc-ch9-2-2-schwarz}. Divide by $b-a>0$.
  $\square$
\end{proof}

\begin{lemma}[the equality case]
  \label{lem:dc-ch9-2-3-equality-case}
  \dcref{ch9:2.3}
  \lean{Riemannian.dcSpeed_ae_const_of_dcEnergy_eq,
    Riemannian.dcArcLength_eq_of_dcEnergy_eq}
  \leanok
  \uses{lem:dc-ch9-2-3-geodesic-equality, lem:dc-ch9-2-2-schwarz, def:dc-ch9-2-2-energy}
  In the situation of \cref{lem:dc-ch9-2-3-inequality}, if moreover $E(\gamma)=E(c)$, then
  \emph{(a)} the speed $|dc/dt|$ is (almost everywhere) constant, i.e. the parameter of $c$
  is proportional to arc length; and \emph{(b)} $L(\gamma)=L(c)$, i.e. $c$ is minimizing too.

  These are the two conclusions do Carmo extracts and feeds to \cref{cor:dc-ch3-3-9} to
  conclude that $c$ is a geodesic. They are \emph{weaker} than that corollary's literal
  hypotheses, which is exactly why that last step is still open: the corollary wants arc
  length proportional to the parameter \emph{pointwise} (not merely almost everywhere),
  and minimality against \emph{every} competing curve (not merely $L(\gamma)=L(c)$ for the
  one $c$ at hand).
\end{lemma}
\begin{proof}
  \leanok
  \uses{lem:dc-ch9-2-3-geodesic-equality, lem:dc-ch9-2-2-schwarz}
  Under $E(\gamma)=E(c)$ every inequality in the chain
  $(b-a)E(\gamma)=L(\gamma)^2\le L(c)^2\le (b-a)E(c)$ of
  \cref{lem:dc-ch9-2-3-inequality} collapses to an equality. From $L(c)^2=(b-a)E(c)$ and the
  equality case of \cref{lem:dc-ch9-2-2-schwarz}, the speed of $c$ is a.e. constant, which is
  (a); from $L(\gamma)^2=L(c)^2$ and non-negativity of both lengths, $L(\gamma)=L(c)$, which
  is (b). $\square$
\end{proof}

Every formula of this section is written in terms of $\frac{D}{dt}$, the covariant
derivative along a curve of Prop. 2.2, Chap. 2 (\cref{prop:dc-ch2-2-2}). That proposition
is available only in a \emph{fixed chart}: it reads the curve and the field in one
coordinate system, so it cannot be stated for a curve leaving that chart --- and the
geodesics of \cref{thm:dc-ch9-3-1} and \cref{thm:dc-ch9-3-7} do leave every chart. The
two nodes below record the chart-independent form of $\frac{D}{dt}$ and of the product
rule of Prop. 3.2, Chap. 2 (\cref{prop:dc-ch2-3-2}) that the formulas of this section
actually use; they carry $\frac{DV}{dt}$ as a second field $DV$ rather than as an applied
operator, which is what makes them chart-local.

\begin{definition}[the covariant derivative along a curve, chart-independently]
  \label{def:dc-ch9-2-covariant-pair}
  \dcref{ch9:2.4}
  \lean{Riemannian.Variation.IsCovariantDerivSolOn,
    Riemannian.Variation.IsCovariantDerivFieldAlongOn,
    Riemannian.Variation.IsCovariantDerivSolOn.transfer,
    Riemannian.Variation.IsCovariantDerivFieldAlongOn.isCovariantDerivSolOn_of_mem_source}
  \leanok
  \uses{prop:dc-ch2-2-2, lem:dc-ch2-2-2-coord}
  Let $c:I\to M$ be a curve and let $V$, $DV$ be vector fields along $c$. Say that
  \emph{$DV$ is the covariant derivative of $V$ along $c$} on $[a,b]$ if near every
  $t_0\in[a,b]$ there is a chart containing the nearby piece of $c$ in which the
  coordinate readings satisfy formula (1) of \cref{lem:dc-ch2-2-2-coord}, namely
  $\dot V=DV-\Gamma(\dot u,V)(u)$, i.e. $\frac{DV}{dt}=DV$.

  The condition does not depend on the charts chosen: under a change of chart the
  inhomogeneous second-derivative term of the transition map produced by the product rule
  cancels against the transformation law of the Christoffel symbols, which is exactly the
  statement that $\frac{D}{dt}$ is a geometric operator. Being chart-local, the notion is
  meaningful for a curve that leaves every single chart, and on any subinterval lying in
  one chart it localizes back to the coordinate equation.

  (Only differentiability of the curve read in charts is required --- not that $c$ be a
  geodesic. The variations of \cref{def:dc-ch9-2-1} have non-geodesic transversals, so a
  geodesic hypothesis here would exclude precisely the fields this section is about.)
\end{definition}

\begin{lemma}[product rule along a curve, chart-independently]
  \label{lem:dc-ch9-2-covariant-pair-leibniz}
  \dcref{ch9:2.4}
  \lean{Riemannian.Variation.IsCovariantDerivFieldAlongOn.hasDerivAt_metricInner}
  \leanok
  \uses{def:dc-ch9-2-covariant-pair, lem:dc-ch2-3-2-coord}
  Let $V$, $W$ be vector fields along a curve $c$ with covariant derivatives $DV$, $DW$
  in the sense of \cref{def:dc-ch9-2-covariant-pair}. Then, at interior times,
  $$
  \frac{d}{dt}\langle V,W\rangle=\Big\langle\frac{DV}{dt},W\Big\rangle
    +\Big\langle V,\frac{DW}{dt}\Big\rangle.
  $$
  This is formula (3) of Prop. 3.2, Chap. 2 in the chart-independent form used throughout
  this section.
\end{lemma}
\begin{proof}
  \leanok
  \uses{def:dc-ch9-2-covariant-pair, lem:dc-ch2-3-2-coord}
  The claim is local in $t$, so fix an interior time and a chart containing the nearby
  piece of $c$, as provided by \cref{def:dc-ch9-2-covariant-pair}. In that chart the
  intrinsic pairing $\langle V,W\rangle$ is the Gram pairing of the coordinate readings,
  and the coordinate metric compatibility of $\frac{D}{dt}$ (\cref{lem:dc-ch2-3-2-coord},
  the Levi-Civita case of formula (3) of \cref{prop:dc-ch2-3-2}) gives its derivative as
  $\langle\frac{DV}{dt},W\rangle+\langle V,\frac{DW}{dt}\rangle$ with $\frac{D}{dt}$ the
  coordinate covariant derivative of \cref{lem:dc-ch2-2-2-coord}. By
  \cref{def:dc-ch9-2-covariant-pair} those two coordinate covariant derivatives are the
  readings of $DV$ and $DW$, which is the assertion. $\square$
\end{proof}

The proof of the first variation formula below exchanges $\frac{D}{ds}\frac{\partial f}{\partial t}$
for $\frac{D}{dt}\frac{\partial f}{\partial s}$ "using the symmetry of the Riemannian
connection" --- the symmetry lemma of Lemma 3.4, Chap. 3 (\cref{lem:dc-ch3-3-4}) read on the
parametrized surface $f$. That lemma is available as an identity between \emph{coordinate
expressions}; the exchange the proof performs is between the two covariant derivative
\emph{operators} $\frac{D}{\partial s}$, $\frac{D}{\partial t}$ along a surface. The node
below records the operator form, which is the shape both \cref{prop:dc-ch9-2-4} and
\cref{prop:dc-ch9-2-8} use.

\begin{lemma}[the symmetry lemma, operator form]
  \label{lem:dc-ch9-2-4-symmetry}
  \dcref{ch9:2.4}
  \lean{Riemannian.Variation.surfaceCovariantDerivS_snd_eq_surfaceCovariantDerivT_fst,
    Riemannian.Variation.surfaceCovariantDerivS_snd_eq_surfaceCovariantDerivT_fst_of_forall}
  \leanok
  \uses{lem:dc-ch3-3-4, def:dc-ch3-3-3}
  Let $f$ be a ($C^2$) parametrized surface (\cref{def:dc-ch3-3-3}) \emph{read in a fixed
  chart}, and let $\frac{D}{\partial s}$, $\frac{D}{\partial t}$ be the covariant
  derivatives along $f$ of that chart. Then the two exchange on the velocity fields of $f$:
  $$
  \frac{D}{\partial s}\frac{\partial f}{\partial t}
    = \frac{D}{\partial t}\frac{\partial f}{\partial s}.
  $$
  No curvature term appears. Contrast \cref{lem:dc-ch4-4-1}, which exchanges
  $\frac{D}{\partial s}$ and $\frac{D}{\partial t}$ on an \emph{arbitrary} field $V$ along
  $f$ and does pick up $R(\frac{\partial f}{\partial s},\frac{\partial f}{\partial t})V$;
  here the two sides differ only by the mixed second partial, which Schwarz kills.

  (The statement is \emph{chart-fixed}, exactly like \cref{lem:dc-ch3-3-4} on which it
  rests and like \cref{lem:dc-ch4-4-1}: $f$ is the reading of the surface in one chart, and
  both operators are that chart's. This suffices for the proofs of \cref{prop:dc-ch9-2-4}
  and \cref{prop:dc-ch9-2-8}, which differentiate under the integral sign and may work
  locally in $t$. A chart-independent form --- of the kind
  \cref{def:dc-ch9-2-covariant-pair} provides for the one-parameter $\frac{D}{dt}$ --- is
  \emph{not} claimed here: no chart-free two-parameter surface exists in the formalization
  yet.)
\end{lemma}
\begin{proof}
  \leanok
  \uses{lem:dc-ch3-3-4}
  Both sides expand, in a chart, as a mixed second partial of $f$ plus a Christoffel term:
  $\frac{D}{\partial s}\frac{\partial f}{\partial t}
   = \frac{\partial^2 f}{\partial s\,\partial t}
     + \Gamma\big(\frac{\partial f}{\partial s},\frac{\partial f}{\partial t}\big)(f)$, and
  symmetrically with $s$, $t$ interchanged. The mixed partials agree by Schwarz, and the two
  Christoffel terms agree because $\Gamma$ is symmetric in its two vector slots --- which is
  exactly the hypothesis that the connection is symmetric. This is \cref{lem:dc-ch3-3-4} at
  $v=(1,0)$, $w=(0,1)$. $\square$
\end{proof}

do Carmo's proof of \cref{prop:dc-ch9-2-4} splits into two halves that are independent of
each other, and the two nodes below record them separately. The \emph{surface half}
(\cref{lem:dc-ch9-2-4-energy-density}) differentiates the energy density in $s$ and applies
the symmetry lemma; it is pointwise, and chart-fixed because
\cref{lem:dc-ch9-2-4-symmetry} is. The \emph{intrinsic half}
(\cref{lem:dc-ch9-2-4-parts}) integrates the result by parts; it mentions no surface at
all --- only two covariant-derivative pairs along $c$ in the sense of
\cref{def:dc-ch9-2-covariant-pair} --- and is therefore chart-independent, so it holds for
a curve leaving every chart.

The surface half admits a \emph{chart-free} formulation, and it is the one the two halves
compose through. The chart in \cref{lem:dc-ch9-2-4-energy-density} is an artefact of the
route, not of the mathematics: the pointwise step is metric compatibility along a
\emph{transversal} $\sigma\mapsto f(\sigma,t)$, and a transversal is a curve like any
other, so \cref{def:dc-ch9-2-covariant-pair} already covers it with no chart and no
two-parameter surface object. \cref{lem:dc-ch9-2-4-density-self} records that form,
\cref{lem:dc-ch9-2-4-under-integral} performs the exchange of $\frac{d}{ds}$ with
$\int dt$, and \cref{lem:dc-ch9-2-4-segment} assembles formula (1) on a segment.

\cref{prop:dc-ch9-2-4} still carries no \lean{} tag, for three reasons. First, its symmetry
hypothesis $\frac{D}{\partial s}\frac{\partial f}{\partial t}=\frac{D}{\partial t}\frac{\partial f}{\partial s}$
is supplied on all of $(a,b)$ at once, whereas \cref{lem:dc-ch9-2-4-symmetry-manifold}
supplies it at one point at a time under a chart condition on the two slice lines through
that point; assembling the pointwise identities into the $(a,b)$-wide hypothesis is a
separate step. Second, the covariant-pair hypotheses of \cref{lem:dc-ch9-2-4-segment} are,
today, only produced along a geodesic (through the Jacobi-field bridge); realizing them for
the transversals of a general variation --- do Carmo's $f(s,t)=\exp_{c(t)}sV(t)$ of
\cref{prop:dc-ch9-2-2} --- is not yet formalized, so the segment lemma has no call site that
starts from a bare variation. Third, the statement is do Carmo's $k=0$ case, so the jump sum
over his subdivision is not yet formed. This is an accounting of what is missing today, not
a claim that nothing else is.

\begin{lemma}[the $s$-derivative of the energy density]
  \label{lem:dc-ch9-2-4-energy-density}
  \dcref{ch9:2.4}
  \lean{Riemannian.Variation.hasDerivAt_chartEnergyDensity_slice,
    Riemannian.Variation.hasDerivAt_chartEnergyDensity_slice_symm}
  \leanok
  \uses{lem:dc-ch9-2-4-symmetry, lem:dc-ch2-3-2-coord}
  Let $f$ be a ($C^2$) parametrized surface \emph{read in a fixed chart}, as in
  \cref{lem:dc-ch9-2-4-symmetry}. Then, pointwise,
  $$
  \frac{\partial}{\partial s}\Big\langle\frac{\partial f}{\partial t},\frac{\partial f}{\partial t}\Big\rangle
    = 2\Big\langle\frac{D}{\partial s}\frac{\partial f}{\partial t},\frac{\partial f}{\partial t}\Big\rangle
    = 2\Big\langle\frac{D}{\partial t}\frac{\partial f}{\partial s},\frac{\partial f}{\partial t}\Big\rangle.
  $$
  This is the step "differentiating under the integral sign and using the symmetry of the
  Riemannian connection" in do Carmo's proof of \cref{prop:dc-ch9-2-4}, with the
  differentiation under the integral sign removed: what remains is the identity the
  integrand satisfies at each point. (The statement is chart-fixed, inherited from
  \cref{lem:dc-ch9-2-4-symmetry}; no chart-free two-parameter surface exists in the
  formalization yet.)
\end{lemma}
\begin{proof}
  \leanok
  \uses{lem:dc-ch9-2-4-symmetry, lem:dc-ch2-3-2-coord}
  The first equality is coordinate metric compatibility (\cref{lem:dc-ch2-3-2-coord})
  applied along the $s$-slice with $V=W=\frac{\partial f}{\partial t}$: the product rule's
  two terms $\langle\frac{D}{\partial s}V,W\rangle+\langle V,\frac{D}{\partial s}W\rangle$
  coincide because the pairing is symmetric, which supplies the factor $2$. The second
  equality is \cref{lem:dc-ch9-2-4-symmetry}, the symmetry of the connection, applied to
  the first factor. No curvature appears. $\square$
\end{proof}

\begin{lemma}[integration by parts along a curve]
  \label{lem:dc-ch9-2-4-parts}
  \dcref{ch9:2.4}
  \lean{Riemannian.Variation.IsCovariantDerivFieldAlongOn.integral_metricInner_add,
    Riemannian.Variation.IsCovariantDerivFieldAlongOn.integral_metricInner_covariantDeriv_left,
    Riemannian.Variation.IsCovariantDerivFieldAlongOn.integral_metricInner_eq_neg_integral_of_proper}
  \leanok
  \uses{lem:dc-ch9-2-covariant-pair-leibniz, def:dc-ch9-2-covariant-pair}
  Let $V$, $W$ be vector fields along a curve $c$ with covariant derivatives $DV$, $DW$
  in the sense of \cref{def:dc-ch9-2-covariant-pair}, on a segment $[a,b]$ carrying no
  breakpoint, and suppose the two pairings $\langle\frac{DV}{dt},W\rangle$ and
  $\langle V,\frac{DW}{dt}\rangle$ are integrable on $[a,b]$ (do Carmo's curves are only
  piecewise differentiable, so this is not automatic and cannot be dropped --- without it
  the displayed identity is not even well posed). Then
  $$
  \int_a^b\Big\langle\frac{DV}{dt},W\Big\rangle dt
    = \langle V(b),W(b)\rangle-\langle V(a),W(a)\rangle
      -\int_a^b\Big\langle V,\frac{DW}{dt}\Big\rangle dt.
  $$
  If moreover $V(a)=V(b)=0$ --- do Carmo's properness condition --- the two boundary terms
  drop, and the identity reads
  $\int_a^b\langle\frac{DV}{dt},W\rangle dt=-\int_a^b\langle V,\frac{DW}{dt}\rangle dt$.

  Taking $W=\frac{dc}{dt}$ and $DW=\frac{D}{dt}\frac{dc}{dt}$, this is the right-hand side
  of formula (1) of \cref{prop:dc-ch9-2-4} on a single segment: the boundary terms are do
  Carmo's $-\langle V(0),\frac{dc}{dt}(0)\rangle+\langle V(a),\frac{dc}{dt}(a)\rangle$, and
  summing over the segments $[t_i,t_{i+1}]$ of his subdivision telescopes the boundary
  terms at the interior breakpoints into his jump sum
  $-\sum_i\langle V(t_i),\frac{dc}{dt}(t_i^+)-\frac{dc}{dt}(t_i^-)\rangle$.
\end{lemma}
\begin{proof}
  \leanok
  \uses{lem:dc-ch9-2-covariant-pair-leibniz}
  By \cref{lem:dc-ch9-2-covariant-pair-leibniz}, $\frac{d}{dt}\langle V,W\rangle
  =\langle\frac{DV}{dt},W\rangle+\langle V,\frac{DW}{dt}\rangle$ at interior times, and
  $t\mapsto\langle V,W\rangle$ is continuous on $[a,b]$. The fundamental theorem of
  calculus (in the form requiring continuity on $[a,b]$ and a derivative on $(a,b)$ only,
  which is what \cref{lem:dc-ch9-2-covariant-pair-leibniz} supplies) integrates this to
  $\int_a^b(\langle\frac{DV}{dt},W\rangle+\langle V,\frac{DW}{dt}\rangle)dt
  =\langle V(b),W(b)\rangle-\langle V(a),W(a)\rangle$; splitting the integral and
  rearranging gives the claim. When $V(a)=V(b)=0$ both boundary terms vanish by linearity
  of the pairing in its first slot. $\square$
\end{proof}

\begin{lemma}[the $s$-derivative of the energy density, chart-free]
  \label{lem:dc-ch9-2-4-density-self}
  \dcref{ch9:2.4}
  \lean{Riemannian.Variation.IsCovariantDerivFieldAlongOn.hasDerivAt_metricInner_self}
  \leanok
  \uses{lem:dc-ch9-2-covariant-pair-leibniz, def:dc-ch9-2-covariant-pair}
  Let $V$ be a vector field along a curve $\gamma$ with covariant derivative $DV$ in the
  sense of \cref{def:dc-ch9-2-covariant-pair}. Then, at interior times,
  $$
  \frac{d}{dt}\langle V,V\rangle = 2\Big\langle\frac{DV}{dt},V\Big\rangle.
  $$

  Applied along a \emph{transversal} $\sigma\mapsto f(\sigma,t)$ of a variation
  (\cref{def:dc-ch9-2-1}) with $V=\frac{\partial f}{\partial t}$, this is the first
  equality of \cref{lem:dc-ch9-2-4-energy-density} with the chart removed: a transversal is
  a curve like any other, and \cref{def:dc-ch9-2-covariant-pair} is chart-independent, so
  no two-parameter surface object is needed to state or to prove it.
\end{lemma}
\begin{proof}
  \leanok
  \uses{lem:dc-ch9-2-covariant-pair-leibniz}
  By \cref{lem:dc-ch9-2-covariant-pair-leibniz} at $W=V$, one has
  $\frac{d}{dt}\langle V,V\rangle
  =\langle\frac{DV}{dt},V\rangle+\langle V,\frac{DV}{dt}\rangle$ at interior times. The two
  summands are equal by the symmetry of the metric, which supplies the factor $2$.
  $\square$
\end{proof}

\begin{lemma}[symmetry of the connection along a surface, chart-free]
  \label{lem:dc-ch9-2-4-symmetry-manifold}
  \dcref{ch9:2.4}
  \lean{Riemannian.Variation.covariantDerivS_velT_eq_covariantDerivT_velS}
  \leanok
  \uses{lem:dc-ch9-2-4-symmetry, def:dc-ch9-2-covariant-pair}
  Let $f$ be a parametrized surface in $M$, let $\frac{\partial f}{\partial t}$ and
  $\frac{\partial f}{\partial s}$ be its coordinate velocity fields, and suppose
  $\big(\frac{\partial f}{\partial t},\frac{D}{\partial s}\frac{\partial f}{\partial t}\big)$
  is a covariant pair along the transversal through $(s_0,t_0)$ and
  $\big(\frac{\partial f}{\partial s},\frac{D}{\partial t}\frac{\partial f}{\partial s}\big)$
  a covariant pair along the curve in the variation through $(s_0,t_0)$, both in the sense
  of \cref{def:dc-ch9-2-covariant-pair}. If $f$ is $C^2$ near $(s_0,t_0)$ and the two slice
  curves through $(s_0,t_0)$ stay in the source of a single chart on their windows, then
  $$
  \frac{D}{\partial s}\frac{\partial f}{\partial t}(s_0,t_0)
    = \frac{D}{\partial t}\frac{\partial f}{\partial s}(s_0,t_0).
  $$

  This is \cref{lem:dc-ch9-2-4-symmetry} with the chart discharged: the conclusion is an
  identity between intrinsic covariant derivatives, and the chart appears only as a device
  in the proof. The hypothesis is on the two \emph{slice lines} through the point, not on
  the surface: no other slice is constrained.
\end{lemma}
\begin{proof}
  \leanok
  \uses{lem:dc-ch9-2-4-symmetry, def:dc-ch9-2-covariant-pair}
  Localize both covariant pairs into the common chart at $\alpha$: on a window whose slice
  curve lies in that chart's source, a covariant pair in the sense of
  \cref{def:dc-ch9-2-covariant-pair} restricts to the coordinate certificate of that chart,
  and at interior times the certificate identifies the intrinsic covariant derivative with
  the coordinate operator $\frac{D}{\partial\cdot}$ applied to the chart reading
  $F=\varphi_\alpha\circ f$. The coordinate operators of the two slices of $F$ are exactly
  the two operators of \cref{lem:dc-ch9-2-4-symmetry}, whose conclusion therefore gives
  the identity between the chart readings of the two sides. The chart reading of an
  intrinsic vector is its image under the tangent coordinate change, which is a linear
  isomorphism; cancelling it gives the identity between the intrinsic vectors themselves.
  $\square$
\end{proof}

\begin{lemma}[differentiation of the energy under the integral sign]
  \label{lem:dc-ch9-2-4-under-integral}
  \dcref{ch9:2.4}
  \lean{Riemannian.Variation.hasDerivAt_dcEnergy_of_dominated}
  \leanok
  \uses{lem:dc-ch9-2-4-density-self, def:dc-ch9-2-2-energy, def:dc-ch9-2-1,
    def:dc-ch9-2-covariant-pair}
  Let $f$ be a variation of $c$ and let $E$ be its energy (\cref{def:dc-ch9-2-2-energy}).
  Suppose $\big(\frac{\partial f}{\partial t},\frac{D}{\partial s}\frac{\partial f}{\partial t}\big)$
  is a covariant pair along each transversal, that the energy density is integrable at
  $s_0$, and that on a neighbourhood of $s_0$ the integrands
  $\big\langle\frac{D}{\partial s}\frac{\partial f}{\partial t},\frac{\partial f}{\partial t}\big\rangle$
  are dominated by a fixed integrable function of $t$. Then
  $$
  E'(s_0)
    = 2\int_a^b\Big\langle\frac{D}{\partial s}\frac{\partial f}{\partial t},\frac{\partial f}{\partial t}\Big\rangle\Big|_{s_0} dt .
  $$

  The domination hypothesis is not removable at this generality, and it is what do Carmo's
  "differentiating under the integral sign" silently assumes. For a $C^2$ variation on a
  compact parameter box it is automatic: the integrand is continuous on a compact set, so a
  constant dominates it.
\end{lemma}
\begin{proof}
  \leanok
  \uses{lem:dc-ch9-2-4-density-self, def:dc-ch9-2-2-energy}
  By \cref{def:dc-ch9-2-2-energy}, $E(s)=\int_a^b\langle\frac{\partial f}{\partial t},\frac{\partial f}{\partial t}\rangle dt$
  with the integrand evaluated along the curve $f(s,\cdot)$. For each fixed $t$ the map
  $\sigma\mapsto\langle\frac{\partial f}{\partial t},\frac{\partial f}{\partial t}\rangle(\sigma,t)$
  is, by \cref{lem:dc-ch9-2-4-density-self} applied along the transversal through $t$,
  differentiable at each interior $\sigma$ with derivative
  $2\langle\frac{D}{\partial s}\frac{\partial f}{\partial t},\frac{\partial f}{\partial t}\rangle(\sigma,t)$.
  Together with integrability at $s_0$ and the domination hypothesis, the standard theorem
  on differentiation under the integral sign applies and yields the displayed identity.
  $\square$
\end{proof}

\begin{lemma}[formula (1) on a segment carrying no breakpoint]
  \label{lem:dc-ch9-2-4-segment}
  \dcref{ch9:2.4}
  \lean{Riemannian.Variation.hasDerivAt_dcEnergy_eq_first_variation}
  \leanok
  \uses{lem:dc-ch9-2-4-under-integral, lem:dc-ch9-2-4-parts, def:dc-ch9-2-2-energy,
    def:dc-ch9-2-covariant-pair}
  In the situation of \cref{lem:dc-ch9-2-4-under-integral}, let $V=\frac{\partial f}{\partial s}(s_0,\cdot)$
  carry the covariant derivative $\frac{DV}{dt}$ and let $\frac{dc}{dt}=\frac{\partial f}{\partial t}(s_0,\cdot)$
  carry $\frac{D}{dt}\frac{dc}{dt}$, both along $c=f(s_0,\cdot)$ and in the sense of
  \cref{def:dc-ch9-2-covariant-pair}. Assume the symmetry of the connection along $c$,
  $$
  \frac{D}{\partial s}\frac{\partial f}{\partial t}(s_0,t)=\frac{D}{\partial t}\frac{\partial f}{\partial s}(s_0,t)
  \qquad\text{for every interior }t\in(a,b),
  $$
  and the integrability of the two pairings of \cref{lem:dc-ch9-2-4-parts}. Then
  $$
  \frac{1}{2}E'(s_0)
    = -\int_a^b\Big\langle V,\frac{D}{dt}\frac{dc}{dt}\Big\rangle dt
      -\Big\langle V(a),\frac{dc}{dt}(a)\Big\rangle
      +\Big\langle V(b),\frac{dc}{dt}(b)\Big\rangle .
  $$

  This is formula (1) of \cref{prop:dc-ch9-2-4} in do Carmo's case $k=0$: a variation of a
  curve differentiable across $[a,b]$, so that the subdivision is trivial and the jump sum
  is empty. The symmetry hypothesis is do Carmo's step "using the symmetry of the
  Riemannian connection"; \cref{lem:dc-ch9-2-4-symmetry-manifold} establishes it at one
  point at a time, under a chart condition on the two slice lines through that point.
\end{lemma}
\begin{proof}
  \leanok
  \uses{lem:dc-ch9-2-4-under-integral, lem:dc-ch9-2-4-parts}
  By \cref{lem:dc-ch9-2-4-under-integral},
  $E'(s_0)=2\int_a^b\langle\frac{D}{\partial s}\frac{\partial f}{\partial t},\frac{\partial f}{\partial t}\rangle dt$.
  The symmetry hypothesis rewrites the integrand as
  $\langle\frac{DV}{dt},\frac{dc}{dt}\rangle$ pointwise, so
  $\frac{1}{2}E'(s_0)=\int_a^b\langle\frac{DV}{dt},\frac{dc}{dt}\rangle dt$. Applying
  \cref{lem:dc-ch9-2-4-parts} with $W=\frac{dc}{dt}$ and $DW=\frac{D}{dt}\frac{dc}{dt}$
  turns the right-hand side into
  $\langle V(b),\frac{dc}{dt}(b)\rangle-\langle V(a),\frac{dc}{dt}(a)\rangle
  -\int_a^b\langle V,\frac{D}{dt}\frac{dc}{dt}\rangle dt$, which is the claim. $\square$
\end{proof}

\begin{proposition}[formula for the first variation of the energy of a curve]
  \label{prop:dc-ch9-2-4}
  \dcref{ch9:2.4}
  \uses{def:dc-ch9-2-1, def:dc-ch9-2-2-energy, lem:dc-ch9-2-4-symmetry,
    def:dc-ch9-2-covariant-pair, lem:dc-ch9-2-4-energy-density, lem:dc-ch9-2-4-parts,
    lem:dc-ch9-2-4-density-self, lem:dc-ch9-2-4-symmetry-manifold,
    lem:dc-ch9-2-4-under-integral, lem:dc-ch9-2-4-segment}
  Let $c:[0,a]\to M$ be a piecewise differentiable curve and let
  $f:(-\varepsilon,\varepsilon)\times[0,a]\to M$ be a variation of $c$. If
  $E:(-\varepsilon,\varepsilon)\to\mathbf{R}$ is the energy of $f$ then

  $$
  \frac{1}{2}E'(0)=-\int_0^a\Big\langle V(t),\frac{D}{dt}\frac{dc}{dt}\Big\rangle dt-\sum_{i=1}^k\Big\langle V(t_i),\frac{dc}{dt}(t_i^+)-\frac{dc}{dt}(t_i^-)\Big\rangle-\Big\langle V(0),\frac{dc}{dt}(0)\Big\rangle+\Big\langle V(a),\frac{dc}{dt}(a)\Big\rangle,\qquad(1)
  $$

  where $V(t)$ is the variational field of $f$, and
  $\frac{dc}{dt}(t_i^+)=\lim_{t\to t_i,\,t>t_i}\frac{dc}{dt}$,
  $\frac{dc}{dt}(t_i^-)=\lim_{t\to t_i,\,t<t_i}\frac{dc}{dt}$.
\end{proposition}
\begin{proof}
  \uses{def:dc-ch9-2-2-energy, lem:dc-ch9-2-4-energy-density, lem:dc-ch9-2-4-parts,
    def:dc-ch9-2-covariant-pair}
  By definition,
  $E(s)=\int_0^a\langle\frac{\partial f}{\partial t},\frac{\partial f}{\partial t}\rangle dt=\sum_{i=0}^k\int_{t_i}^{t_{i+1}}\langle\frac{\partial f}{\partial t},\frac{\partial f}{\partial t}\rangle dt$.
  Differentiating under the integral sign and using the symmetry of the Riemannian
  connection,

  $$
  \frac{d}{ds}\int_{t_i}^{t_{i+1}}\Big\langle\frac{\partial f}{\partial t},\frac{\partial f}{\partial t}\Big\rangle dt=2\Big\langle\frac{\partial f}{\partial s},\frac{\partial f}{\partial t}\Big\rangle\Big|_{t_i}^{t_{i+1}}-2\int_{t_i}^{t_{i+1}}\Big\langle\frac{\partial f}{\partial s},\frac{D}{dt}\frac{\partial f}{\partial t}\Big\rangle dt.
  $$

  Therefore,

  $$
  \frac{1}{2}\frac{dE}{ds}=\sum_{i=0}^k\Big\langle\frac{\partial f}{\partial s},\frac{\partial f}{\partial t}\Big\rangle\Big|_{t_i}^{t_{i+1}}-\int_0^a\Big\langle\frac{\partial f}{\partial s},\frac{D}{dt}\frac{\partial f}{\partial t}\Big\rangle dt.\qquad(2)
  $$

  Putting $s=0$ in (2) yields (1). $\square$
\end{proof}

\begin{lemma}[a geodesic is a critical point of the energy]
  \label{lem:dc-ch9-2-5-geodesic}
  \dcref{ch9:2.5}
  \lean{Riemannian.Variation.IsCovariantDerivFieldAlongOn.integral_metricInner_covariantDeriv_eq_zero_of_geodesic}
  \leanok
  \uses{lem:dc-ch9-2-4-parts, def:dc-ch9-2-covariant-pair}
  Let $W$ be a \emph{parallel} field along a curve $c$ --- $\frac{DW}{dt}=0$ --- and let
  $V$, $DV$ be a covariant-derivative pair along $c$ with $V(a)=V(b)=0$. Then on the
  segment $[a,b]$
  $$
  \int_a^b\Big\langle\frac{DV}{dt},W\Big\rangle dt=0.
  $$

  Taking $W=\frac{dc}{dt}$ --- so that the hypothesis $\frac{DW}{dt}=0$ becomes
  $\frac{D}{dt}\frac{dc}{dt}=0$, i.e.\ $c$ is a geodesic --- and taking $V$ to be the
  variational field of a proper variation, so that $V(a)=V(b)=0$ is do Carmo's properness
  condition, this is the direction \emph{geodesic $\Rightarrow$ critical point} of
  \cref{prop:dc-ch9-2-5}: in do Carmo's words, "all terms of (1) are zero" --- the
  integral term because $\frac{D}{dt}\frac{dc}{dt}=0$, the boundary terms because the
  variation is proper. Reading that vanishing integral \emph{as} $\frac12\frac{dE}{ds}(0)$
  is the content of \cref{prop:dc-ch9-2-4}, which is not yet proved; the statement above
  --- the part that is formalized --- asserts only the vanishing of the integral, and says
  nothing about $E$.
\end{lemma}
\begin{proof}
  \leanok
  \uses{lem:dc-ch9-2-4-parts}
  Apply \cref{lem:dc-ch9-2-4-parts} with $DW=0$. The remaining integral
  $\int_a^b\langle V,\frac{DW}{dt}\rangle dt$ vanishes because its integrand is
  identically zero, and the two boundary terms vanish because $V(a)=V(b)=0$. $\square$
\end{proof}

\begin{proposition}[geodesics as critical points of energy]
  \label{prop:dc-ch9-2-5}
  \dcref{ch9:2.5}
  \uses{prop:dc-ch9-2-4, prop:dc-ch9-2-2, lem:dc-ch9-2-5-geodesic}
  A piecewise differentiable curve $c:[0,a]\to M$ is a geodesic if and only if, for
  every proper variation $f$ of $c$, we have $\frac{dE}{ds}(0)=0$.
\end{proposition}
\begin{proof}
  \uses{prop:dc-ch9-2-4, prop:dc-ch9-2-2, lem:dc-ch9-2-5-geodesic}
  If $c$ is a geodesic, $\frac{D}{dt}\frac{dc}{dt}=0$ and $c$ is regular;
  for $f$ proper, $V(0)=V(a)=0$, and all terms of (1) are zero
  (\cref{lem:dc-ch9-2-5-geodesic}). Conversely, suppose
  $\frac{dE}{ds}(0)=0$ for all proper variations. Let
  $V(t)=g(t)\frac{D}{dt}\frac{dc}{dt}$, where $g:[0,a]\to\mathbf{R}$ is piecewise
  differentiable with $g(t)>0$ if $t\ne t_i$ and $g(t_i)=0$, $i=0,\dots,k+1$.
  Constructing a variation
  with this $V(t)$,

  $$
  \frac{1}{2}\frac{dE}{ds}(0)=-\int_0^a g(t)\Big\langle\frac{D}{dt}\frac{dc}{dt},\frac{D}{dt}\frac{dc}{dt}\Big\rangle dt=0,
  $$

  so $\frac{D}{dt}\frac{dc}{dt}=0$ on each $(t_i,t_{i+1})$, that is, $c$ is a geodesic
  on each $(t_i,t_{i+1})$. To handle the points $t_i$, take another variational field
  $\overline V(t)$ with $\overline V(0)=\overline V(a)=0$ and, for $t\ne 0,a$,
  $\overline V(t_i)=\frac{dc}{dt}(t_i^+)-\frac{dc}{dt}(t_i^-)$. Then

  $$
  0=\frac{1}{2}\frac{dE}{ds}(0)=-\sum_{i=1}^k\Big|\frac{dc}{dt}(t_i^+)-\frac{dc}{dt}(t_i^-)\Big|^2,
  $$

  so $c$ is of class $C^1$ at each $t_i$. Since $\frac{D}{dt}\frac{dc}{dt}=0$ at
  $t_i$, $c$ satisfies the geodesic equation on $(0,a)$; by uniqueness of solutions
  to ordinary differential equations, $c\in C^\infty$ and is a geodesic. $\square$
\end{proof}

\begin{remark}[Remark 2.6]
  \label{rem:dc-ch9-2-6}
  \dcref{ch9:2.6}
  \uses{prop:dc-ch9-2-5}
  \cref{prop:dc-ch9-2-5} furnishes a characterization of geodesics as critical points of
  the energy for all proper variations. In this sense geodesics can be thought of as
  solutions to a variational problem; such a characterization is not local but
  involves the behavior of the geodesic as a whole.
\end{remark}

\begin{remark}[Remark 2.7]
  \label{rem:dc-ch9-2-7}
  \dcref{ch9:2.7}
  There is an analogy between the first variation of a proper variation and the usual
  derivative of a function on a differentiable manifold. It is natural to think of
  the set $\Omega_{p,q}$ of piecewise differentiable curves $c$ joining two points
  $p$ and $q$ of $M$ as a manifold, a tangent vector at the point $c$ being a
  piecewise differentiable vector field $V$ along $c$ which vanishes at the
  extremities of $c$. The energy $E$ is then a differentiable function on such a
  manifold, and $\frac{dE}{ds}(0)$ is the derivative of $E$ in the direction of $V$;
  geodesics joining $p$ to $q$ are critical points of $E$. The tangent space at a
  point $c$ does not have finite dimension, so local parametrizations by open subsets
  of $\mathbf{R}^n$ are impossible; this suggests manifolds with infinitely many
  dimensions (consult Palais, R., "Morse theory on Hilbert manifolds", Topology 2
  (1963), 299–340).
\end{remark}

do Carmo's proof of \cref{prop:dc-ch9-2-8} "takes the derivative of (2)". Just as for the
first variation, that splits into an \emph{analytic} step --- differentiating the first
variation $E'(s)=2\int\langle\frac{D}{\partial s}\frac{\partial f}{\partial t},\frac{\partial f}{\partial t}\rangle\,dt$
of \cref{lem:dc-ch9-2-4-under-integral} once more in $s$ --- and a \emph{geometric} step ---
substituting the symmetry of the connection (\cref{lem:dc-ch9-2-4-symmetry}) and the Ricci
identity (\cref{lem:dc-ch4-4-1}) to turn the resulting integrand into the index form
$\langle V',V'\rangle-\langle R(\gamma',V)\gamma',V\rangle$ at $s=0$ for a geodesic. The two
nodes below record the analytic step, which is chart-free;
\cref{lem:dc-ch9-2-8-curvature-substitution} records the geometric substitution --- the Ricci
identity read metric-paired against the intrinsic curvature form --- at the chart-read level,
and \cref{lem:dc-ch9-2-8-curvature-substitution-manifold} lifts it chart-free.
\cref{lem:dc-ch9-2-8-segment} then \emph{assembles} the two into formula (3) on a single
segment carrying no breakpoint, for a proper variation of a geodesic --- do Carmo's $k=0$
case. The general piecewise \cref{prop:dc-ch9-2-8}, with its jump sum over a subdivision, is
obtained by summing that segment identity and is not yet performed, so \cref{prop:dc-ch9-2-8}
itself carries no \lean{} tag.

\begin{lemma}[the second differentiation under the integral sign]
  \label{lem:dc-ch9-2-8-under-integral}
  \dcref{ch9:2.8}
  \lean{Riemannian.Variation.hasDerivAt_dcPairing_of_dominated}
  \leanok
  \uses{lem:dc-ch9-2-covariant-pair-leibniz, def:dc-ch9-2-covariant-pair,
    lem:dc-ch9-2-4-under-integral}
  Let $f$ be a variation of $c$, let $\frac{\partial f}{\partial t}$ carry the covariant
  $s$-derivative $\frac{D}{\partial s}\frac{\partial f}{\partial t}$ along each transversal,
  and let $\frac{D}{\partial s}\frac{\partial f}{\partial t}$ in turn carry its own covariant
  $s$-derivative $\frac{D}{\partial s}\frac{D}{\partial s}\frac{\partial f}{\partial t}$, all
  in the sense of \cref{def:dc-ch9-2-covariant-pair}; suppose the second integrand is
  dominated near $s_0$ by a fixed integrable function of $t$. Then
  $$
  \frac{d}{ds}\int_a^b\Big\langle\frac{D}{\partial s}\frac{\partial f}{\partial t},
    \frac{\partial f}{\partial t}\Big\rangle dt\Big|_{s_0}
  = \int_a^b\Big\{\Big\langle\frac{D}{\partial s}\frac{D}{\partial s}\frac{\partial f}{\partial t},
    \frac{\partial f}{\partial t}\Big\rangle
    + \Big\langle\frac{D}{\partial s}\frac{\partial f}{\partial t},
      \frac{D}{\partial s}\frac{\partial f}{\partial t}\Big\rangle\Big\}\,dt .
  $$
  The pointwise input is metric compatibility along each transversal
  (\cref{lem:dc-ch9-2-covariant-pair-leibniz}) applied to the pairs
  $\big(\frac{D}{\partial s}\frac{\partial f}{\partial t},
    \frac{D}{\partial s}\frac{D}{\partial s}\frac{\partial f}{\partial t}\big)$ and
  $\big(\frac{\partial f}{\partial t},\frac{D}{\partial s}\frac{\partial f}{\partial t}\big)$;
  the exchange of $\frac{d}{ds}$ with $\int dt$ is the same differentiation-under-the-integral
  theorem used in \cref{lem:dc-ch9-2-4-under-integral}.
\end{lemma}
\begin{proof}
  \leanok
  \uses{lem:dc-ch9-2-covariant-pair-leibniz}
  For each fixed $t$ the map
  $\sigma\mapsto\langle\frac{D}{\partial s}\frac{\partial f}{\partial t},\frac{\partial f}{\partial t}\rangle(\sigma,t)$
  is differentiable at interior $\sigma$ with derivative
  $\langle\frac{D}{\partial s}\frac{D}{\partial s}\frac{\partial f}{\partial t},\frac{\partial f}{\partial t}\rangle
  +\langle\frac{D}{\partial s}\frac{\partial f}{\partial t},\frac{D}{\partial s}\frac{\partial f}{\partial t}\rangle$,
  by \cref{lem:dc-ch9-2-covariant-pair-leibniz} (the product rule, not its $V=W$ collapse).
  Together with the domination hypothesis the standard theorem on differentiation under the
  integral sign yields the displayed identity. $\square$
\end{proof}

\begin{lemma}[the second variation of the energy, before the curvature substitution]
  \label{lem:dc-ch9-2-8-energy-snd-deriv}
  \dcref{ch9:2.8}
  \lean{Riemannian.Variation.hasDerivAt_deriv_dcEnergy_second_variation}
  \leanok
  \uses{lem:dc-ch9-2-8-under-integral, lem:dc-ch9-2-4-under-integral,
    def:dc-ch9-2-2-energy, def:dc-ch9-2-covariant-pair}
  In the situation of \cref{lem:dc-ch9-2-8-under-integral}, suppose moreover that the first
  variation $E'(\sigma)=2\int_a^b\langle\frac{D}{\partial s}\frac{\partial f}{\partial t},
  \frac{\partial f}{\partial t}\rangle\,dt$ of \cref{lem:dc-ch9-2-4-under-integral} holds at
  every $\sigma$ in a neighbourhood of $s_0$. Then
  $$
  \tfrac12 E''(s_0)
  = \int_a^b\Big\{\Big\langle\frac{D}{\partial s}\frac{D}{\partial s}\frac{\partial f}{\partial t},
    \frac{\partial f}{\partial t}\Big\rangle
    + \Big\langle\frac{D}{\partial s}\frac{\partial f}{\partial t},
      \frac{D}{\partial s}\frac{\partial f}{\partial t}\Big\rangle\Big\}\,dt .
  $$

  This is do Carmo's step "taking the derivative of (2)", with the \emph{curvature}
  substitution not yet performed: on a neighbourhood of $s_0$ the derivative $E'$ \emph{is}
  $2\int\langle\frac{D}{\partial s}\frac{\partial f}{\partial t},\frac{\partial f}{\partial t}\rangle\,dt$,
  so $E''$ is the derivative of that expression, supplied by
  \cref{lem:dc-ch9-2-8-under-integral}. At $s=0$ for a geodesic, the symmetry of the
  connection and the Ricci identity (\cref{lem:dc-ch4-4-1}) rewrite the integrand as
  $\langle V',V'\rangle-\langle R(\gamma',V)\gamma',V\rangle$ --- the index form
  $I_a(V,V)$ of \cref{def:dc-ch9-2-10-index-form} --- which is formula (6) of
  \cref{rem:dc-ch9-2-10}; that geometric step is a separate node, still open.
\end{lemma}
\begin{proof}
  \leanok
  \uses{lem:dc-ch9-2-8-under-integral, lem:dc-ch9-2-4-under-integral}
  The first variation, supplied on a neighbourhood of $s_0$, says that $E'$ agrees there with
  $\sigma\mapsto 2\int_a^b\langle\frac{D}{\partial s}\frac{\partial f}{\partial t},\frac{\partial f}{\partial t}\rangle\,dt$;
  hence $E''(s_0)$ is the derivative of the latter at $s_0$, which
  \cref{lem:dc-ch9-2-8-under-integral} computes. Transferring the derivative along the
  eventual equality gives the claim. $\square$
\end{proof}

\begin{lemma}[the Ricci-identity substitution, metric-paired, chart-read]
  \label{lem:dc-ch9-2-8-curvature-substitution}
  \dcref{ch9:2.8}
  \lean{Riemannian.Variation.chartMetricInner_surface_covariant_commutator_eq_curvatureFormAt,
    Riemannian.Variation.chartMetricInner_chartCurvatureContraction2_eq_curvatureFormAt,
    Riemannian.Variation.chartMetricInner_chartCurvatureContraction2_chartVectorRep_eq_curvatureFormAt}
  \leanok
  \uses{lem:dc-ch4-4-1, def:dc-ch4-2-1}
  Let $f$ be a ($C^2$) parametrized surface \emph{read in a fixed chart} and $V$ a field along
  it. The commutator of the two covariant derivatives, paired against a vector $w$, is the
  intrinsic curvature $(0,4)$ form of the chart-frame realizations of $\frac{\partial f}{\partial s}$,
  $\frac{\partial f}{\partial t}$, $V$ and $w$:
  $$
  \Big\langle\frac{D}{\partial t}\frac{D}{\partial s}V - \frac{D}{\partial s}\frac{D}{\partial t}V,\ w\Big\rangle
    = \Big\langle R\Big(\frac{\partial f}{\partial s},\frac{\partial f}{\partial t}\Big)V,\ w\Big\rangle,
  $$
  in do Carmo's Ch. 4 Def. 2.1 curvature convention. This is Lemma 4.1 of Chap. 4
  (\cref{lem:dc-ch4-4-1}) read metric-paired and carried across to the intrinsic curvature form
  (\cref{def:dc-ch4-2-1}); it is the geometric substitution do Carmo performs at the curvature
  step of \cref{prop:dc-ch9-2-8}.

  (Chart-fixed, exactly like \cref{lem:dc-ch9-2-4-symmetry}: the covariant derivatives are the
  chart surface operators and the four vectors survive as chart-frame realizations. The fully
  chart-free form --- with the covariant derivatives as intrinsic
  \cref{def:dc-ch9-2-covariant-pair} pairs, of the kind
  \cref{lem:dc-ch9-2-4-symmetry-manifold} provides for the symmetry lemma --- is not claimed
  here.)
\end{lemma}
\begin{proof}
  \leanok
  \uses{lem:dc-ch4-4-1}
  The chart-level Ricci identity (\cref{lem:dc-ch4-4-1}) equates the commutator with the
  coordinate curvature contraction $\mathcal{R}_{\mathrm{chart}}(\frac{\partial f}{\partial s},\frac{\partial f}{\partial t})V$.
  Pairing against $w$ under the chart Gram form and applying the manifold--chart curvature bridge
  --- with antisymmetry in the first pair cancelling the do Carmo -- Morgan--Tian sign carried by
  that bridge --- rewrites it as the intrinsic curvature $(0,4)$ form
  $R(\frac{\partial f}{\partial s},\frac{\partial f}{\partial t},V,w)$. $\square$
\end{proof}

\begin{lemma}[the Ricci-identity substitution, chart-free]
  \label{lem:dc-ch9-2-8-curvature-substitution-manifold}
  \dcref{ch9:2.8}
  \lean{Riemannian.Variation.metricInner_covariantDerivT_covariantDerivS_commutator_eq_curvatureFormAt}
  \leanok
  \uses{lem:dc-ch9-2-8-curvature-substitution, def:dc-ch9-2-covariant-pair,
    lem:dc-ch9-2-4-symmetry-manifold}
  Let $f$ be a parametrized surface in $M$ and $V$ a field along it, and suppose
  $\frac{D}{\partial s}V$, $\frac{D}{\partial t}V$ are covariant derivatives of $V$ along the
  slices in the sense of \cref{def:dc-ch9-2-covariant-pair}, and
  $\frac{D}{\partial t}\frac{D}{\partial s}V$, $\frac{D}{\partial s}\frac{D}{\partial t}V$ the
  covariant derivatives of those, along the two slice curves through $(s_0,t_0)$. Then, at
  $(s_0,t_0)$,
  $$
  \Big\langle\frac{D}{\partial t}\frac{D}{\partial s}V - \frac{D}{\partial s}\frac{D}{\partial t}V,\ W\Big\rangle
    = \Big\langle R\Big(\frac{\partial f}{\partial s},\frac{\partial f}{\partial t}\Big)V,\ W\Big\rangle
  $$
  for every $W\in T_{f(s_0,t_0)}M$, with $R$ the curvature of Def. 2.1, Chap. 4.

  This is \cref{lem:dc-ch9-2-8-curvature-substitution} with the chart discharged: the
  conclusion is an identity between intrinsic covariant derivatives, the covariant
  derivatives are \cref{def:dc-ch9-2-covariant-pair} pairs, and the chart appears only as a
  device in the proof --- the exact analogue, for the Ricci identity, of what
  \cref{lem:dc-ch9-2-4-symmetry-manifold} provides for the symmetry lemma. The hypotheses
  constrain $\frac{D}{\partial s}V$, $\frac{D}{\partial t}V$ on the two-dimensional window
  (the outer covariant derivatives differentiate them along the transverse slice), and
  $\frac{D}{\partial t}\frac{D}{\partial s}V$, $\frac{D}{\partial s}\frac{D}{\partial t}V$ only
  on their own slice line.
\end{lemma}
\begin{proof}
  \leanok
  \uses{lem:dc-ch9-2-8-curvature-substitution, lem:dc-ch9-2-4-symmetry-manifold}
  Localize each of the four covariant pairs into a common chart at $\alpha$ (as in
  \cref{lem:dc-ch9-2-4-symmetry-manifold}), so that at interior times each covariant
  derivative is its chart surface operator applied to the chart readings; the inner
  identification is an equality of fields near the point, so it survives the outer
  differentiation. The chart-read Ricci identity
  (\cref{lem:dc-ch9-2-8-curvature-substitution}) then applies to the two iterated operators,
  and the readback of the chart-frame realizations, together with the velocity bridge for
  $\frac{\partial f}{\partial s}$, $\frac{\partial f}{\partial t}$, collapses the curvature
  side to the intrinsic vectors. $\square$
\end{proof}

\begin{lemma}[formula (3) on a segment, for a proper variation]
  \label{lem:dc-ch9-2-8-segment}
  \dcref{ch9:2.8}
  \lean{Riemannian.Variation.deriv_deriv_dcEnergy_eq_second_variation,
    Riemannian.Variation.integral_second_variation_integrand_eq,
    Riemannian.Variation.curvatureFormAt_swap_pairs}
  \leanok
  \uses{lem:dc-ch9-2-8-energy-snd-deriv, lem:dc-ch9-2-4-symmetry-manifold,
    lem:dc-ch9-2-8-curvature-substitution-manifold, lem:dc-ch9-2-4-parts,
    lem:dc-ch9-2-5-geodesic, def:dc-ch9-2-2-energy, def:dc-ch9-2-covariant-pair}
  Let $\gamma=f(s_0,\cdot):[a,b]\to M$ be a geodesic and $f$ a variation of $\gamma$ with
  variational field $V=\frac{\partial f}{\partial s}(s_0,\cdot)$ satisfying $V(a)=V(b)=0$ and
  $\frac{D}{\partial s}\frac{\partial f}{\partial s}(s_0,a)
    =\frac{D}{\partial s}\frac{\partial f}{\partial s}(s_0,b)=0$ (all of which hold when $f$ is
  proper). Then, on the segment $[a,b]$ carrying no breakpoint,
  $$
  \frac{1}{2}E''(s_0)=-\int_a^b\Big\langle V,\frac{D^2 V}{dt^2}+R(\gamma',V)\gamma'\Big\rangle dt .
  $$

  This is formula (3) of \cref{prop:dc-ch9-2-8} in do Carmo's case $k=0$: the subdivision is
  trivial, so the jump sum is empty, and the variation being proper kills every boundary term.
\end{lemma}
\begin{proof}
  \leanok
  \uses{lem:dc-ch9-2-8-energy-snd-deriv, lem:dc-ch9-2-4-symmetry-manifold,
    lem:dc-ch9-2-8-curvature-substitution-manifold, lem:dc-ch9-2-4-parts, lem:dc-ch9-2-5-geodesic}
  By \cref{lem:dc-ch9-2-8-energy-snd-deriv},
  $\frac12 E''(s_0)=\int_a^b\{\langle\frac{D}{\partial s}\frac{D}{\partial s}\frac{\partial f}{\partial t},\frac{\partial f}{\partial t}\rangle+\langle\frac{D}{\partial s}\frac{\partial f}{\partial t},\frac{D}{\partial s}\frac{\partial f}{\partial t}\rangle\}\,dt$.
  The symmetry of the connection (\cref{lem:dc-ch9-2-4-symmetry-manifold}) rewrites the second
  integrand as $\langle V',V'\rangle$, which integration by parts
  (\cref{lem:dc-ch9-2-4-parts}, the variation being proper) turns into
  $-\int_a^b\langle V,\frac{D^2V}{dt^2}\rangle\,dt$. For the first integrand, the symmetry of
  the connection followed by the Ricci identity
  (\cref{lem:dc-ch9-2-8-curvature-substitution-manifold}) gives
  $\langle\frac{D}{\partial s}\frac{D}{\partial s}\frac{\partial f}{\partial t},\frac{\partial f}{\partial t}\rangle
    =\langle\frac{D}{\partial t}\frac{D}{\partial s}\frac{\partial f}{\partial s},\frac{\partial f}{\partial t}\rangle
     -R\big(\tfrac{\partial f}{\partial s},\tfrac{\partial f}{\partial t},\tfrac{\partial f}{\partial s},\tfrac{\partial f}{\partial t}\big)$,
  whose first term integrates to $0$ because $\gamma$ is a geodesic and
  $\frac{D}{\partial s}\frac{\partial f}{\partial s}$ vanishes at the endpoints
  (\cref{lem:dc-ch9-2-5-geodesic}), and whose curvature term is $-\langle R(\gamma',V)\gamma',V\rangle$:
  at $s=0$, $\frac{\partial f}{\partial s}=V$ and $\frac{\partial f}{\partial t}=\gamma'$, so it reads
  $R(V,\gamma',V,\gamma')$, which the interchange of the two antisymmetric pairs turns into
  $R(\gamma',V,\gamma',V)=\langle R(\gamma',V)\gamma',V\rangle$.
  Summing gives formula (3). $\square$
\end{proof}

\begin{proposition}[formula for the second variation]
  \label{prop:dc-ch9-2-8}
  \dcref{ch9:2.8}
  Let $\gamma:[0,a]\to M$ be a geodesic and let
  $f:(-\varepsilon,\varepsilon)\times[0,a]\to M$ be a proper variation of $\gamma$. Let $E$
  be the energy function of the variation. Then

  $$
  \frac{1}{2}E''(0)=-\int_0^a\Big\langle V(t),\frac{D^2 V}{dt^2}+R\Big(\frac{d\gamma}{dt},V\Big)\frac{d\gamma}{dt}\Big\rangle dt-\sum_{i=1}^k\Big\langle V(t_i),\frac{DV}{dt}(t_i^+)-\frac{DV}{dt}(t_i^-)\Big\rangle,\qquad(3)
  $$

  where $V$ is the variational field of $f$, $R$ is the curvature of $M$ and
  $\frac{DV}{dt}(t_i^+)=\lim_{t\to t_i,\,t>t_i}\frac{DV}{dt}$,
  $\frac{DV}{dt}(t_i^-)=\lim_{t\to t_i,\,t<t_i}\frac{DV}{dt}$.
\end{proposition}
\begin{proof}
  Taking the derivative of (2),

  $$
  \frac{1}{2}\frac{d^2 E}{ds^2}=\sum_{i=0}^k\Big\langle\frac{D}{ds}\frac{\partial f}{\partial s},\frac{\partial f}{\partial t}\Big\rangle\Big|_{t_i}^{t_{i+1}}+\sum_{i=0}^k\Big\langle\frac{\partial f}{\partial s},\frac{D}{ds}\frac{\partial f}{\partial t}\Big\rangle\Big|_{t_i}^{t_{i+1}}-\int_0^a\Big\langle\frac{D}{ds}\frac{\partial f}{\partial s},\frac{D}{dt}\frac{\partial f}{\partial t}\Big\rangle dt-\int_0^a\Big\langle\frac{\partial f}{\partial s},\frac{D}{ds}\frac{D}{dt}\frac{\partial f}{\partial t}\Big\rangle dt.
  $$

  Putting $s=0$, the first and third terms are zero, since $f$ is proper and
  $\gamma$ is a geodesic. Because

  $$
  \frac{D}{ds}\frac{D}{dt}\frac{\partial f}{\partial t}=\frac{D}{dt}\frac{D}{ds}\frac{\partial f}{\partial t}+R\Big(\frac{\partial f}{\partial t},\frac{\partial f}{\partial s}\Big)\frac{\partial f}{\partial t},
  $$

  we have, at $s=0$,
  $\frac{D}{ds}\frac{D}{dt}\frac{\partial f}{\partial t}=\frac{D^2}{dt^2}V+R(\frac{d\gamma}{dt},V)\frac{d\gamma}{dt}$.
  Further use of the fact that the variation is proper yields

  $$
  \sum_{i=0}^k\Big\langle\frac{\partial f}{\partial s},\frac{D}{ds}\frac{\partial f}{\partial t}\Big\rangle\Big|_{t_i}^{t_{i+1}}=-\sum_{i=1}^k\Big\langle V(t_i),\frac{DV}{dt}(t_i^+)-\frac{DV}{dt}(t_i^-)\Big\rangle.\qquad(4)
  $$

  Putting all these facts together, we obtain (3). $\square$
\end{proof}

\begin{remark}[Remark 2.9]
  \label{rem:dc-ch9-2-9}
  \dcref{ch9:2.9}
  \uses{prop:dc-ch9-2-8}
  If the variation is not proper, the first term at the beginning of \cref{prop:dc-ch9-2-8}
  need not be zero. Taking this with the corresponding terms at $i=0$ and $i=k+1$ of
  the second member of (4), we obtain

  $$
  \frac{1}{2}E''(0)=-\int_0^a\Big\langle V(t),\frac{D^2 V}{dt^2}+R\Big(\frac{d\gamma}{dt},V\Big)\frac{d\gamma}{dt}\Big\rangle dt-\sum_{i=1}^k\Big\langle V(t_i),\frac{DV}{dt}(t_i^+)-\frac{DV}{dt}(t_i^-)\Big\rangle
  $$

  $$
  -\Big\langle\frac{D}{ds}\frac{\partial f}{\partial s},\frac{d\gamma}{dt}\Big\rangle(0,0)+\Big\langle\frac{D}{ds}\frac{\partial f}{\partial s},\frac{d\gamma}{dt}\Big\rangle(0,a)-\Big\langle V(0),\frac{DV}{dt}(0)\Big\rangle+\Big\langle V(a),\frac{DV}{dt}(a)\Big\rangle.\qquad(5)
  $$
\end{remark}

\begin{lemma}[formula (5) on a segment, for a non-proper variation]
  \label{lem:dc-ch9-2-9-segment}
  \dcref{ch9:2.9}
  \lean{Riemannian.Variation.deriv_deriv_dcEnergy_eq_second_variation_nonproper,
    Riemannian.Variation.integral_second_variation_integrand_eq_nonproper}
  \leanok
  \uses{lem:dc-ch9-2-8-energy-snd-deriv, lem:dc-ch9-2-4-symmetry-manifold,
    lem:dc-ch9-2-8-curvature-substitution-manifold, lem:dc-ch9-2-4-parts,
    def:dc-ch9-2-2-energy, def:dc-ch9-2-covariant-pair}
  Let $\gamma=f(s_0,\cdot):[a,b]\to M$ be a geodesic and $f$ a variation of $\gamma$ that need
  \emph{not} be proper. Then, on the segment $[a,b]$ carrying no breakpoint,
  $$
  \frac{1}{2}E''(s_0)=-\int_a^b\Big\langle V,\frac{D^2 V}{dt^2}+R(\gamma',V)\gamma'\Big\rangle dt
    +\Big\langle\frac{D}{\partial s}\frac{\partial f}{\partial s},\gamma'\Big\rangle\Big|_a^b
    +\big\langle V,V'\big\rangle\Big|_a^b .
  $$

  This is formula (5) of \cref{rem:dc-ch9-2-9} in do Carmo's case $k=0$ (empty jump sum): the
  two integrations by parts keep their boundary terms instead of dropping them, giving exactly
  do Carmo's extra terms at $i=0$ and $i=k+1$. It is the form \cref{thm:dc-ch9-3-7}
  (Synge--Weinstein) applies, whose variation $h(s,t)=\exp_{\gamma(t)}(s\,e_1(t))$ is not proper.
\end{lemma}
\begin{proof}
  \leanok
  \uses{lem:dc-ch9-2-8-energy-snd-deriv, lem:dc-ch9-2-4-symmetry-manifold,
    lem:dc-ch9-2-8-curvature-substitution-manifold, lem:dc-ch9-2-4-parts}
  Identical to the assembly of \cref{lem:dc-ch9-2-8-segment}, except that the two integrations
  by parts (\cref{lem:dc-ch9-2-4-parts}) are applied in their boundary-keeping form: the
  $\langle V',V'\rangle$ term integrates to
  $\langle V,V'\rangle|_a^b-\int_a^b\langle V,\frac{D^2V}{dt^2}\rangle$, and the geodesic term
  $\langle\frac{D}{\partial t}\frac{D}{\partial s}\frac{\partial f}{\partial s},\gamma'\rangle$
  integrates to $\langle\frac{D}{\partial s}\frac{\partial f}{\partial s},\gamma'\rangle|_a^b$
  (the remaining integral vanishes as $\gamma$ is a geodesic). The curvature term is unchanged
  from \cref{lem:dc-ch9-2-8-segment}. $\square$
\end{proof}

\begin{remark}[the index form]
  \label{rem:dc-ch9-2-10}
  \dcref{ch9:2.10}
  \uses{prop:dc-ch9-2-8, rem:dc-ch9-2-9}
  It is often convenient to write (5) using, on each interval where $V$ is
  differentiable,
  $\frac{d}{dt}\langle V,\frac{DV}{dt}\rangle=\langle V,\frac{D^2 V}{dt^2}\rangle+\langle\frac{DV}{dt},\frac{DV}{dt}\rangle$.
  Taking a geodesic $\gamma:[0,a]\to M$ and a partition
  $0=t_0<t_1<\cdots<t_{k+1}=a$, this gives

  $$
  \frac{1}{2}E''(0)=\int_0^a\{\langle V',V'\rangle-\langle R(\gamma',V)\gamma',V\rangle\}dt-\Big\langle\frac{D}{ds}\frac{\partial f}{\partial s},\gamma'\Big\rangle(0,0)+\Big\langle\frac{D}{ds}\frac{\partial f}{\partial s},\gamma'\Big\rangle(0,a).
  $$

  For reasons which will be clear later, it is convenient to write

  $$
  \int_0^a\{\langle V',V'\rangle-\langle R(\gamma',V)\gamma',V\rangle\}dt=I_a(V,V).
  $$

  If the variation is proper, $I_a(V,V)=\frac{1}{2}E''(0)$ depends only on $V$. In
  the general case,

  $$
  \frac{1}{2}E''(0)=I_a(V,V)-\Big\langle\frac{D}{ds}\frac{\partial f}{\partial s},\gamma'\Big\rangle(0,0)+\Big\langle\frac{D}{ds}\frac{\partial f}{\partial s},\gamma'\Big\rangle(0,a).\qquad(6)
  $$
\end{remark}

The index form $I_a$ is named here because it is studied in its own right later, and the
differentiation identity the remark opens with is the step that turns (5) into (6); both
are therefore recorded as separate nodes.

\begin{lemma}[differentiating $\langle V,DV/dt\rangle$]
  \label{lem:dc-ch9-2-10-differentiation}
  \dcref{ch9:2.10}
  \lean{Riemannian.Variation.hasDerivAt_metricInner_covariantDeriv}
  \leanok
  \uses{lem:dc-ch9-2-covariant-pair-leibniz}
  On each interval where $V$ is differentiable,
  $$
  \frac{d}{dt}\Big\langle V,\frac{DV}{dt}\Big\rangle
    =\Big\langle\frac{DV}{dt},\frac{DV}{dt}\Big\rangle
     +\Big\langle V,\frac{D^2V}{dt^2}\Big\rangle,
  $$
  the identity used to pass from (5) to (6).
\end{lemma}
\begin{proof}
  \leanok
  \uses{lem:dc-ch9-2-covariant-pair-leibniz}
  Apply the product rule \cref{lem:dc-ch9-2-covariant-pair-leibniz} to the pair $V$, with
  covariant derivative $\frac{DV}{dt}$, against the pair $\frac{DV}{dt}$, with covariant
  derivative $\frac{D^2V}{dt^2}$. $\square$
\end{proof}

\begin{definition}[the index form $I_a$]
  \label{def:dc-ch9-2-10-index-form}
  \dcref{ch9:2.10}
  \lean{Riemannian.Variation.indexForm, Riemannian.Variation.indexForm_def}
  \leanok
  \uses{def:dc-ch9-2-covariant-pair, def:dc-ch4-2-1, def:dc-ch1-2-9}
  For a geodesic $\gamma:[0,a]\to M$ and a vector field $V$ along $\gamma$ with covariant
  derivative $V'=\frac{DV}{dt}$, the \emph{index form} is
  $$
  I_a(V,V)=\int_0^a\{\langle V',V'\rangle-\langle R(\gamma',V)\gamma',V\rangle\}\,dt,
  $$
  with $R$ the curvature of Def. 2.1, Chap. 4 and $\gamma'$ the velocity of Def. 2.9,
  Chap. 1. This is a definition only: the identity $I_a(V,V)=\frac{1}{2}E''(0)$ for a
  proper variation is formula (6) above, recorded (on a single segment) as
  \cref{lem:dc-ch9-2-10-formula6}.
\end{definition}

\begin{lemma}[formula (6): the second variation is the index form, on a segment]
  \label{lem:dc-ch9-2-10-formula6}
  \dcref{ch9:2.10}
  \lean{Riemannian.Variation.deriv_deriv_dcEnergy_eq_indexForm,
    Riemannian.Variation.integral_second_variation_eq_indexForm}
  \leanok
  \uses{def:dc-ch9-2-10-index-form, lem:dc-ch9-2-8-energy-snd-deriv,
    lem:dc-ch9-2-4-symmetry-manifold, lem:dc-ch9-2-8-curvature-substitution-manifold,
    lem:dc-ch9-2-5-geodesic}
  Let $\gamma=f(s_0,\cdot):[a,b]\to M$ be a geodesic and $f$ a proper variation of $\gamma$
  (in the sense of \cref{lem:dc-ch9-2-8-segment}). Then
  $$
  \frac{1}{2}E''(0)=I_a(V,V).
  $$
  This is formula (6) of \cref{rem:dc-ch9-2-10} for a proper variation on a single segment ---
  the same assembly as \cref{lem:dc-ch9-2-8-segment}, read one integration by parts earlier, so
  that the $\langle V',V'\rangle$ term is kept rather than pushed back to
  $-\langle V,\frac{D^2 V}{dt^2}\rangle$.
\end{lemma}
\begin{proof}
  \leanok
  \uses{lem:dc-ch9-2-8-energy-snd-deriv, lem:dc-ch9-2-4-symmetry-manifold,
    lem:dc-ch9-2-8-curvature-substitution-manifold, lem:dc-ch9-2-5-geodesic,
    def:dc-ch9-2-10-index-form}
  By \cref{lem:dc-ch9-2-8-energy-snd-deriv},
  $\frac12 E''(0)=\int_a^b\{\langle\frac{D}{\partial s}\frac{D}{\partial s}\frac{\partial f}{\partial t},\frac{\partial f}{\partial t}\rangle+\langle\frac{D}{\partial s}\frac{\partial f}{\partial t},\frac{D}{\partial s}\frac{\partial f}{\partial t}\rangle\}\,dt$.
  The symmetry of the connection (\cref{lem:dc-ch9-2-4-symmetry-manifold}) makes the second
  integrand $\langle V',V'\rangle$; the symmetry followed by the Ricci identity
  (\cref{lem:dc-ch9-2-8-curvature-substitution-manifold}), together with the geodesic
  vanishing of the boundary term (\cref{lem:dc-ch9-2-5-geodesic}), makes the first integrand
  integrate to $-\int_a^b\langle R(\gamma',V)\gamma',V\rangle\,dt$. Their sum is $I_a(V,V)$.
  $\square$
\end{proof}

\begin{remark}[Remark 2.11]
  \label{rem:dc-ch9-2-11}
  \dcref{ch9:2.11}
  It is possible to establish formulas for the first and second variations of the
  arc length function $L(s)$. The expressions are entirely analogous to those for the
  energy and, since they will not be used here, are not treated (see [dC 2]
  paragraph 5.4).
\end{remark}

\begin{remark}[Remark 2.12]
  \label{rem:dc-ch9-2-12}
  \dcref{ch9:2.12}
  \uses{rem:dc-ch9-2-7}
  The analogy mentioned in \cref{rem:dc-ch9-2-7} can be carried further, considering
  $2I_a(V,V)$ as the hessian $d^2 E(V,V)$ of the energy function
  $E:\Omega_{p,q}\to\mathbf{R}$ at the critical point $\gamma\in\Omega_{p,q}$ with
  respect to the "vector" $V$.
\end{remark}

\section{The theorems of Bonnet–Myers and of Synge–Weinstein}

We now go into some applications of the formula for the second variation of the
energy.

Both applications open the same way: along the geodesic $\gamma$ one chooses \emph{parallel}
orthonormal fields $e_1,\dots,e_{n-1}$ orthogonal to $\gamma'$, together with $e_n=\gamma'/|\gamma'|$.
That frame is the node below; it is the manifold (chart-crossing) counterpart of the chart-level
velocity-seeded frame of Ch.~8, needed here because the Bonnet--Myers geodesic (of length $>\pi r$)
leaves every chart.

\begin{lemma}[a velocity-seeded parallel orthonormal frame along a geodesic]
  \label{lem:dc-ch9-3-1-velocity-frame}
  \dcref{ch9:3.1}
  \lean{Riemannian.Variation.exists_velocitySeededParallelOrthoFrameAlongOn}
  \leanok
  Let $\gamma:[a,b]\to M$ be a geodesic with $\gamma'(a)\ne 0$. Then there are parallel fields
  $e_1,\dots,e_n$ along $\gamma$, orthonormal at every $t\in[a,b]$, and a distinguished index
  $n_0$ with $e_{n_0}(t)=\ell^{-1}\gamma'(t)$ for all $t$, where $\ell=|\gamma'|>0$ is the
  (constant) speed. In particular $e_i(t)\perp\gamma'(t)$ for $i\ne n_0$.
\end{lemma}
\begin{proof}
  \leanok
  Extend the unit velocity $\ell^{-1}\gamma'(a)$ to a $g$-orthonormal basis of $T_{\gamma(a)}M$
  (Householder reflection: \texttt{exists\_metricOrthonormalBasis\_containing\_unit}), and
  parallel-transport it along $\gamma$; parallel transport is a linear isometry, so the frame
  stays orthonormal. The distinguished transported member and $\ell^{-1}\gamma'$ are both parallel
  fields with the same value at $a$, so by uniqueness of parallel transport they agree on
  $[a,b]$. $\square$
\end{proof}

\begin{lemma}[the index form of a scaled parallel field]
  \label{lem:dc-ch9-3-1-index-form-scaled}
  \dcref{ch9:3.1}
  \lean{Riemannian.Variation.indexForm_smul_eq}
  \leanok
  \uses{def:dc-ch9-2-10-index-form, def:dc-ch9-2-covariant-pair}
  For a field $V(t)=\varphi(t)\,e(t)$ with covariant derivative $V'=\varphi'\,e$ (as when $e$ is
  parallel), the index form is
  $$I_a(V,V)=\int_a^b\Big((\varphi')^2\,\langle e,e\rangle
    -\varphi^2\,\langle R(\gamma',e)\gamma',e\rangle\Big)\,dt.$$
  A pointwise multilinear rewrite of $I_a$: the scalars $\varphi',\varphi$ come out of the metric
  term and out of the two $e$-slots of the curvature term. This is do Carmo's integrand for
  $V_j=(\sin\pi t)e_j$ in the proof of \cref{thm:dc-ch9-3-1}.
\end{lemma}
\begin{proof}
  \leanok
  Pull $\varphi'$ out of $\langle\varphi' e,\varphi' e\rangle$ (bilinearity of the metric) and
  $\varphi$ out of each $e$-argument of $\langle R(\gamma',\varphi e)\gamma',\varphi e\rangle$
  (homogeneity of the curvature form in slots 2 and 4). $\square$
\end{proof}

\begin{lemma}[the index form of $V_j$ in the velocity-seeded frame]
  \label{lem:dc-ch9-3-1-index-form-frame}
  \dcref{ch9:3.1}
  \lean{Riemannian.Variation.indexForm_smul_frame_eq}
  \leanok
  \uses{lem:dc-ch9-3-1-index-form-scaled, lem:dc-ch9-3-1-velocity-frame}
  With the frame of \cref{lem:dc-ch9-3-1-velocity-frame} (so $e=e_j$ is a unit field and
  $\gamma'=\ell\,e_n$), for $V=\varphi\,e_j$
  $$I_a(V,V)=\int_a^b\Big((\varphi')^2-\varphi^2\,\ell^2\,\langle R(e_n,e_j)e_n,e_j\rangle\Big)\,dt.$$
  The $\langle e_j,e_j\rangle$ factor is $1$ (orthonormality) and $\gamma'=\ell\,e_n$ pulls
  $\ell^2$ out of the two velocity slots of the curvature term. For $\varphi=\sin\pi t$ this is
  do Carmo's integrand $\sin^2\pi t\,(\pi^2-\ell^2 K(e_n,e_j))$ of \cref{thm:dc-ch9-3-1}
  (using $\langle R(e_n,e_j)e_n,e_j\rangle=K(e_n,e_j)$ for orthonormal $e_n\perp e_j$).
\end{lemma}
\begin{proof}
  \leanok
  \uses{lem:dc-ch9-3-1-index-form-scaled}
  Substitute $\langle e_j,e_j\rangle=1$ and $\gamma'=\ell\,e_n$ into
  \cref{lem:dc-ch9-3-1-index-form-scaled} and use homogeneity of the curvature form in slots
  1 and 3. $\square$
\end{proof}

\begin{lemma}[second-order necessary condition at a minimum]
  \label{lem:dc-ch9-3-1-second-order-min}
  \dcref{ch9:3.1}
  \lean{Riemannian.Variation.isLocalMin_deriv_deriv_nonneg}
  \leanok
  If $f:\mathbf{R}\to\mathbf{R}$ has a local minimum at $x$, is continuous there and is twice
  differentiable at $x$, then $f''(x)\ge 0$. This is the step by which a \emph{minimizing}
  geodesic forces $E_j''(0)\ge 0$ in \cref{thm:dc-ch9-3-1}: each curve of the variation joins
  the same endpoints, so $E_j(s)\ge E_j(0)$, i.e. $s=0$ is a minimum of $E_j$.
\end{lemma}
\begin{proof}
  \leanok
  By contradiction: if $f''(x)<0$ then (by the sufficient converse
  \texttt{isLocalMax\_of\_deriv\_deriv\_neg}) $x$ is a local maximum too, so $f$ is locally
  constant, its derivative is eventually $0$, and by uniqueness of the derivative $f''(x)=0$,
  a contradiction. $\square$
\end{proof}

Summing do Carmo's per-$j$ integrand over $j\ne n_0$, feeding in the Ricci bound, and reading
off a negative summand is one self-contained step, recorded below. It is stated at the level of
the \emph{index form} $I_a$ rather than $E_j''(0)$: the two agree by formula (6)
(\cref{lem:dc-ch9-2-10-formula6}) only through the concrete exponential variation of
\cref{prop:dc-ch9-2-2}, which is not yet formalized, so the node carries exactly the part that
is --- the summed negativity, the frame, and the Ricci sum.

\begin{lemma}[the summed index form is negative]
  \label{lem:dc-ch9-3-1-index-sum-neg}
  \dcref{ch9:3.1}
  \lean{Riemannian.Variation.integral_frame_sum_lt_zero,
    Riemannian.Variation.sum_indexForm_smul_frame_neg,
    Riemannian.Variation.exists_indexForm_smul_frame_neg}
  \leanok
  \uses{lem:dc-ch9-3-1-velocity-frame, lem:dc-ch9-3-1-index-form-frame,
    lem:dc-ch9-3-3-ricci-sectional, def:dc-ch9-2-10-index-form}
  Let $\gamma:[0,1]\to M$ be a geodesic of constant speed $\ell$ with the velocity-seeded
  parallel orthonormal frame of \cref{lem:dc-ch9-3-1-velocity-frame} (so each $e_j$, $j\ne n_0$,
  is a unit field and $\gamma'=\ell\,e_{n_0}$), and put $V_j(t)=(\sin\pi t)\,e_j(t)$. If $n\ge 2$
  (there is at least one index $j\ne n_0$), $r>0$, $\pi r<\ell$, and the raw curvature sum satisfies
  $\sum_{j\ne n_0}\langle R(e_{n_0},e_j)e_{n_0},e_j\rangle\ge (n-1)/r^2$ at every $t\in[0,1]$
  (each term continuous in $t$), then
  $$
  \sum_{j\ne n_0} I_a(V_j,V_j)<0,\qquad\text{hence}\quad I_a(V_j,V_j)<0\ \text{for some }j\ne n_0.
  $$

  This is do Carmo's contradiction step. By \cref{lem:dc-ch9-3-1-index-form-frame} each summand
  is $\int_0^1(\pi^2\cos^2\pi t-\ell^2\sin^2\pi t\,\langle R(e_{n_0},e_j)e_{n_0},e_j\rangle)\,dt$,
  and the raw sum $\sum_{j\ne n_0}\langle R(e_{n_0},e_j)e_{n_0},e_j\rangle$ is the unnormalized
  Ricci form $Q(e_{n_0},e_{n_0})$ (\cref{lem:dc-ch9-3-3-ricci-sectional}), i.e. do Carmo's
  $(n-1)\mathrm{Ric}$; the bound is thus his $\mathrm{Ric}\ge 1/r^2$. do Carmo writes this step
  as $\tfrac12\sum_j E_j''(0)<0$; that reading needs formula (6)
  (\cref{lem:dc-ch9-2-10-formula6}, $\tfrac12 E_j''(0)=I_a(V_j,V_j)$) applied to the concrete
  proper variation with variational field $V_j$ (\cref{prop:dc-ch9-2-2}), still open.
\end{lemma}
\begin{proof}
  \leanok
  \uses{lem:dc-ch9-3-1-index-form-frame}
  For each $j\ne n_0$, \cref{lem:dc-ch9-3-1-index-form-frame} with $\varphi=\sin\pi t$
  (so $\varphi'=\pi\cos\pi t$) gives
  $I_a(V_j,V_j)=\int_0^1\big((\pi\cos\pi t)^2-\sin^2\pi t\,(\ell^2\langle R(e_{n_0},e_j)e_{n_0},e_j\rangle)\big)\,dt$.
  The integrands are continuous, so the finite sum commutes with the integral, and summing over
  the $n-1$ indices $j\ne n_0$ collapses it to
  $$
  \int_0^1\Big((n-1)(\pi\cos\pi t)^2-\sin^2\pi t\,\big(\ell^2\textstyle\sum_{j\ne n_0}\langle R(e_{n_0},e_j)e_{n_0},e_j\rangle\big)\Big)\,dt.
  $$
  Write $S(t)=\ell^2\sum_{j\ne n_0}\langle R(e_{n_0},e_j)e_{n_0},e_j\rangle$. The integrand equals
  $(n-1)\pi^2\cos 2\pi t-\sin^2\pi t\,(S(t)-(n-1)\pi^2)$; the first part integrates to $0$
  (as $\int_0^1\cos 2\pi t\,dt=0$), and the second is the integral of a function that is $\ge 0$
  on $[0,1]$ and strictly positive on $(0,1)$ --- there $\sin\pi t>0$ and, since $\pi r<\ell$
  gives $\pi^2 r^2<\ell^2$, $S(t)\ge\ell^2(n-1)/r^2>(n-1)\pi^2$. Hence the whole integral is
  $<0$. A strictly negative finite sum has a strictly negative summand. $\square$
\end{proof}

\begin{theorem}[Bonnet–Myers]
  \label{thm:dc-ch9-3-1}
  \dcref{ch9:3.1}
  \notready
  \uses{lem:dc-ch9-2-3, prop:dc-ch9-2-8, lem:dc-ch9-3-1-velocity-frame,
    lem:dc-ch9-3-1-index-form-frame, lem:dc-ch9-3-1-index-sum-neg,
    lem:dc-ch9-3-1-second-order-min, lem:dc-ch9-3-3-ricci-sectional}
  Let $M^n$ be a complete Riemannian manifold. Suppose that the Ricci curvature of
  $M$ satisfies

  $$
  \mathrm{Ric}_p(v)\ge\frac{1}{r^2}>0,
  $$

  for all $p\in M$ and for all $v\in T_p(M)$. Then $M$ is compact and the diameter
  $\mathrm{diam}(M)\le\pi r$.
\end{theorem}
\begin{proof}
  Let $p$ and $q$ be any two points in $M$. Since $M$ is complete, there
  exists a minimizing geodesic $\gamma:[0,1]\to M$ joining $p$ to $q$. It is enough
  to show that the length $\ell(\gamma)\le\pi r$, because then $M$ is bounded and
  complete, therefore compact; and since $d(p,q)\le\pi r$ for any $p,q$, it follows
  that $\mathrm{diam}(M)\le\pi r$. Suppose, to the contrary, that
  $\ell(\gamma)=\ell>\pi r$. Consider parallel fields $e_1(t),\dots,e_{n-1}(t)$ along
  $\gamma$ which are orthonormal for each $t\in[0,1]$ and belong to the orthogonal
  complement of $\gamma'(t)$. Let $e_n(t)=\frac{\gamma'(t)}{\ell}$ and let $V_j$ be
  the vector field along $\gamma$ given by $V_j(t)=(\sin\pi t)e_j(t)$,
  $j=1,\dots,n-1$. Then $V_j(0)=V_j(1)=0$, so $V_j$ generates a proper variation of
  $\gamma$ with energy $E_j$. Using the formula for the second variation and the fact
  that $e_j$ is parallel,

  $$
  \frac{1}{2}E_j''(0)=-\int_0^1\langle V_j,V_j''+R(\gamma',V_j)\gamma'\rangle dt=\int_0^1\sin^2\pi t\,(\pi^2-\ell^2 K(e_n(t),e_j(t)))\,dt,
  $$

  where $K(e_n(t),e_j(t))$ is the sectional curvature at $\gamma(t)$ with respect to
  the plane generated by $e_n(t),e_j(t)$. Summing on $j$ and using the definition of
  Ricci curvature,

  $$
  \frac{1}{2}\sum_{j=1}^{n-1}E_j''(0)=\int_0^1\{\sin^2\pi t\,((n-1)\pi^2-(n-1)\ell^2\,\mathrm{Ric}_{\gamma(t)}(e_n(t)))\}\,dt.
  $$

  Since $\mathrm{Ric}_{\gamma(t)}(e_n(t))\ge\frac{1}{r^2}$ and $\ell>\pi r$, we have
  $(n-1)\ell^2\,\mathrm{Ric}_{\gamma(t)}(e_n(t))>(n-1)\pi^2$, hence

  $$
  \frac{1}{2}\sum_{j=1}^{n-1}E_j''(0)<\int_0^1\sin^2\pi t\,((n-1)\pi^2-(n-1)\pi^2)\,dt=0.
  $$

  As a result there exists an index $j$ such that $E_j''(0)<0$, which, by \cref{lem:dc-ch9-2-3},
  contradicts the fact that $\gamma$ is minimizing. Therefore $\ell\le\pi r$.
  $\square$

  In the corollaries that follow, we use some facts on the fundamental group and on
  covering spaces (see Massey [Ma], Chapters 2 and 5).
\end{proof}

\begin{corollary}[finiteness of the fundamental group]
  \label{cor:dc-ch9-3-2}
  \dcref{ch9:3.2}
  Let $M$ be a complete Riemannian manifold with $\mathrm{Ric}_p(v)\ge\delta>0$, for
  all $p\in M$ and all $v\in T_p(M)$. Then the universal cover of $M$ is compact. In
  particular, the fundamental group $\pi_1(M)$ is finite.
\end{corollary}
\begin{proof}
  Introducing on the universal cover $\pi:\tilde M\to M$ the covering metric
  (the metric such that $\pi$ is a local isometry), $\tilde M$ is complete and its
  Ricci curvature satisfies $\widetilde{\mathrm{Ric}}_p\ge\delta>0$. From the theorem,
  $\tilde M$ is compact. Hence the number of sheets of the covering is finite; since
  that is the number of elements in $\pi_1(M)$, we conclude that $\pi_1(M)$ is
  finite. $\square$
\end{proof}

\cref{cor:dc-ch9-3-3} deduces a Ricci hypothesis from a sectional one, a step do Carmo
takes without comment: the Ricci form is a sum of $n-1$ sectional curvatures, so a bound
on every sectional curvature bounds it below. That step is the node below.

\begin{lemma}[Ricci as a sum of sectional curvatures]
  \label{lem:dc-ch9-3-3-ricci-sectional}
  \dcref{ch9:3.3}
  \lean{Riemannian.sectionalCurvature_eq_of_orthonormal,
    Riemannian.ricciForm_self_eq_sum_sectionalCurvature,
    Riemannian.ricciForm_self_ge_of_sectionalCurvature_ge}
  \leanok
  \uses{def:dc-ch4-3-2, def:dc-ch4-4-ricci-form}
  Let $\{e_1,\dots,e_n\}$ be an orthonormal basis of $T_pM$ and put $v=e_{i_0}$. Then
  $$
  Q(v,v)=\sum_{i\ne i_0}K(v,e_i),
  $$
  where $Q$ is the (unnormalized) Ricci form and $K$ the sectional curvature.
  Consequently, if $K(v,e_i)\ge k$ for every $i\ne i_0$, then $Q(v,v)\ge(n-1)k$.

  In do Carmo's normalization $\mathrm{Ric}_p(v)=Q(v,v)/(n-1)$ the consequence reads
  $\mathrm{Ric}_p(v)\ge k$, which is literally the hypothesis of
  \cref{thm:dc-ch9-3-1}; the formalization keeps $Q$ unnormalized, so the factor $n-1$ is
  explicit. The sum runs over $n-1$ terms because the diagonal contribution $B(v,v,v,v)$
  vanishes, and each pair $(v,e_i)$, $i\ne i_0$, being orthonormal, has $|v\wedge e_i|^2=1$,
  so its sectional curvature is the bare numerator $B(v,e_i,v,e_i)$.
\end{lemma}
\begin{proof}
  \leanok
  \uses{def:dc-ch4-3-2, def:dc-ch4-4-ricci-form}
  By the orthonormal-basis formula for the Ricci form, $Q(v,v)=\sum_i B(v,e_i,v,e_i)$. The
  term $i=i_0$ is $B(v,v,v,v)=0$ by the antisymmetry of $B$ in its first two slots. For
  $i\ne i_0$ the pair $(v,e_i)$ is orthonormal, so
  $|v\wedge e_i|^2=\langle v,v\rangle\langle e_i,e_i\rangle-\langle v,e_i\rangle^2=1$ and
  hence $K(v,e_i)=B(v,e_i,v,e_i)/1=B(v,e_i,v,e_i)$. Summing gives the identity. The
  inequality follows by bounding each of the $n-1$ summands below by $k$. $\square$
\end{proof}

\begin{corollary}[positive sectional curvature bound]
  \label{cor:dc-ch9-3-3}
  \dcref{ch9:3.3}
  \uses{thm:dc-ch9-3-1, cor:dc-ch9-3-2, lem:dc-ch9-3-3-ricci-sectional}
  Let $M$ be a complete Riemannian manifold with sectional curvature
  $K\ge\frac{1}{r^2}>0$. Then $M$ is compact, $\mathrm{diam}(M)\le\pi r$ and
  $\pi_1(M)$ is finite.
\end{corollary}

\begin{remark}[Remark 3.4]
  \label{rem:dc-ch9-3-4}
  \dcref{ch9:3.4}
  The hypothesis $K\ge\delta>0$ cannot be relaxed to $K>0$. Indeed, the paraboloid
  $\{(x,y,z)\in\mathbf{R}^3;z=x^2+y^2\}$ has curvature $K>0$, but is complete and
  non-compact.
\end{remark}

\begin{remark}[Remark 3.5]
  \label{rem:dc-ch9-3-5}
  \dcref{ch9:3.5}
  In fact it is not necessary that $K$ be bounded away from zero but only that $K$
  not approach zero too fast. In this respect, see E. Calabi, "On Ricci curvatures
  and geodesics", Duke Math. J. 34 (1967), 667–676 and R. Schneider, "Konvexe
  Flächen mit langsam abnehmender Krümmung", Archiv der Math. 23 (1972), 650–654.
\end{remark}

\begin{remark}[Remark 3.6]
  \label{rem:dc-ch9-3-6}
  \dcref{ch9:3.6}
  \uses{thm:dc-ch9-3-1}
  The estimate for the diameter given by \cref{thm:dc-ch9-3-1} cannot be improved. Indeed, the
  unit sphere $S^n\subset\mathbf{R}^{n+1}$ has constant sectional curvature equal to
  1 (therefore its Ricci curvature also is a constant equal to 1) and
  $\mathrm{diam}(S^n)=\pi$. This example is unique in the following sense: let $M^n$
  be complete with $\mathrm{Ric}_p(v)\ge 1/r^2$, for all $p\in M$ and all $v\in T_pM$;
  if $\mathrm{diam}(M)=\pi r$, then $M^n$ is isometric to the sphere $S^n$ of
  curvature $1/r^2$ (see S.Y. Cheng, "Eigenvalues comparison theorems and its
  geometric applications", Math. Z. 143 (1975), 289–297 and, for another proof, K.
  Shiohama, "A sphere theorem for manifolds of positive Ricci curvature", Trans.
  A.M.S. 275 (1983), 811–819).
\end{remark}

\begin{theorem}[Synge–Weinstein]
  \label{thm:dc-ch9-3-7}
  \dcref{ch9:3.7}
  Let $f$ be an isometry of a compact oriented Riemannian manifold $M^n$. Suppose
  that $M$ has positive sectional curvature and that $f$ preserves the orientation of
  $M$ if $n$ is even, and reverses it if $n$ is odd. Then $f$ has a fixed point, i.e.,
  there exists $p\in M$ with $f(p)=p$.

  In the proof of Theorem 3.7 we need the following fact from linear algebra.
\end{theorem}

\begin{lemma}[invariant vector of an orthogonal map]
  \label{lem:dc-ch9-3-8}
  \dcref{ch9:3.8}
  \lean{Riemannian.Variation.exists_fixed_vector_of_det_eq}
  \leanok
  Let $A$ be an orthogonal linear transformation of $\mathbf{R}^{n-1}$ and suppose
  that $\det A=(-1)^n$. Then $A$ leaves invariant some non-zero vector of
  $\mathbf{R}^{n-1}$.

  (The Lean statement is the same fact for an arbitrary finite-dimensional real inner
  product space $V$ in the role of $\mathbf{R}^{n-1}$, with $n$ pinned by
  $\dim V+1=n$; this is the form in which \cref{thm:dc-ch9-3-7} applies it, namely to
  the orthogonal complement of $\gamma'(0)$ in $T_pM$. This is a statement of pure
  linear algebra: it depends on nothing else in the chapter, and \cref{thm:dc-ch9-3-7}
  depends on \emph{it}, not conversely.)
\end{lemma}
\begin{proof}
  \leanok
  If $n$ is even, $\det(A-\lambda I)$ is a real polynomial in $\lambda$ of
  odd degree $(=n-1)$, so $A$ has a real eigenvalue. Since $A$ is orthogonal, such
  eigenvalues are of the form $\pm 1$. Since the product of the complex eigenvalues
  of $A$ is non-negative, and $\det A=1$, at least one of the eigenvalues equals 1.
  If $n$ is odd, $\det A=-1$. Because the product of the complex eigenvalues is
  non-negative, there is at least one pair of real eigenvalues, one of which is
  positive, hence equal to 1. $\square$

  (The Lean proof establishes the same conclusion by a uniform determinant identity
  rather than by do Carmo's split on the parity of $n$: from $A-I=A(I-A^{T})$ one gets
  $\det(A-I)=\det A\cdot(-1)^{n-1}\cdot\det(A-I)$, and the hypothesis makes that scalar
  $-1$, so $\det(A-I)=0$.)
\end{proof}

\begin{proof}[Proof of \cref{thm:dc-ch9-3-7}]
  \proves{thm:dc-ch9-3-7}
  \uses{rem:dc-ch9-2-9, lem:dc-ch9-3-8, lem:dc-ch9-2-2-schwarz}
  Suppose, to the contrary, that $f(q)\ne q$ for all $q\in M$.
  Let $p\in M$ such that $d(p,f(p))$ attains a minimum. Since $M$ is complete, there
  exists a normalized minimizing geodesic $\gamma:[0,\ell]\to M$, joining $p$ to
  $f(p)$, that is, $\gamma(0)=p$, $\gamma(\ell)=f(p)$. Let
  $\tilde A=P\circ df_p:T_p(M)\to T_p(M)$, where $P$ is the parallel transport along
  $\gamma$ from $f(p)=\gamma(\ell)$ to $p$. Then $\tilde A$ is an isometry. We show
  that $(f\circ\gamma)'(0)=\gamma'(\ell)$. Consider the geodesic $f\circ\gamma$ which
  joins $f(p)$ to $f^2(p)$. Let $p'=\gamma(t')$, $t'\ne 0$, $t'\ne\ell$, and
  $f(p')=f\circ\gamma(t')$. Since $f$ is an isometry, $d(p,p')=d(f(p),f(p'))$ and,
  from the triangle inequality,

  $$
  d(p',f(p'))\le d(p',f(p))+d(f(p),f(p'))=d(p',f(p))+d(p,p')=d(p,f(p)).
  $$

  Because $d(p,f(p))$ is a minimum, $d(p',f(p'))=d(p',f(p))+d(f(p),f(p'))$. Therefore
  the curve formed by $\gamma$ and $f\circ\gamma$ is a geodesic, hence
  $(f\circ\gamma)'(0)=\gamma'(\ell)$, as asserted. It follows that $\tilde A$ leaves
  $\gamma'(0)$ fixed, since

  $$
  \tilde A(\gamma'(0))=P\circ df_p(\gamma'(0))=P((f\circ\gamma)'(0))=P(\gamma'(\ell))=\gamma'(0).
  $$

  Let $A$ be the restriction of $\tilde A$ to the orthogonal complement of
  $\gamma'(0)$. Then $A$ is orthogonal on $\mathbf{R}^{n-1}$ and, since $P$ is an
  orientation-preserving isometry,

  $$
  \det A=\det\tilde A=\det(P\circ df_p)=(-1)^n,
  $$

  where in the last equality we use the hypothesis on $f$ and the fact that $P$
  preserves orientation. From the lemma, $A$ leaves a vector invariant. Let $e_1(t)$
  be a unit parallel field along $\gamma$ such that, for each $t$, $e_1(t)$ belongs to
  the orthogonal complement of $\gamma'(t)$ and $e_1(0)$ is invariant by $A$. Let
  $\beta(s)$, $s\in(-\varepsilon,\varepsilon)$, be a geodesic such that $\beta(0)=p$ and
  $\beta'(0)=e_1(0)$. Because $P\circ df_p(e_1(0))=e_1(0)$, we have
  $df_p(e_1(0))=e_1(\ell)$, that is, the geodesic $f\circ\beta$ is such that
  $f\circ\beta(0)=f(p)$ and $(f\circ\beta)'(0)=e_1(\ell)$. Let $h$ be a variation of
  $\gamma$ given by $h(s,t)=\exp_{\gamma(t)}(se_1(t))$,
  $s\in(-\varepsilon,\varepsilon)$, $t\in[0,\ell]$. Since $h(s,0)=\beta(s)$, then
  $h(s,\ell)=\exp_{f(p)}(se_1(\ell))=(f\circ\beta)(s)$. Therefore,

  $$
  V(t)=\frac{\partial}{\partial s}\exp_{\gamma(t)}(se_1(t))\Big|_{s=0}=e_1(t),
  $$

  hence $\frac{D^2 V}{dt^2}=0$. Using the formula for the second variation (see
  \cref{rem:dc-ch9-2-9}) and the fact that $\frac{\partial h}{\partial s}(0,t)=e_1(t)$,

  $$
  \frac{1}{2}\frac{d^2 E}{ds^2}(0)=-\int_0^\ell\Big\langle V(t),R\Big(\frac{d\gamma}{dt},V\Big)\frac{d\gamma}{dt}\Big\rangle dt+\Big\langle\frac{d\gamma}{dt},\frac{D}{ds}\frac{\partial h}{\partial s}\Big\rangle(0,\ell)-\Big\langle\frac{d\gamma}{dt},\frac{D}{ds}\frac{\partial h}{\partial s}\Big\rangle(0,0)=-\int_0^\ell K\Big(e_1(t),\frac{d\gamma}{dt}\Big)dt.
  $$

  Because the sectional curvature is positive, $\frac{1}{2}\frac{d^2 E}{ds^2}(0)<0$,
  and therefore $\frac{dE}{ds}$ is strictly decreasing on a neighborhood of zero.
  This shows that there exists a curve $c$ in the variation such that
  $(\ell(c))^2\le\ell E(c)<\ell E(\gamma)=\ell(\gamma)^2$. Since the curves in the
  variation join $q$ to $f(q)$, we obtain a contradiction to the fact that
  $d(p,f(p))$ is a minimum. $\square$
\end{proof}

\begin{remark}[Remark 3.9]
  \label{rem:dc-ch9-3-9}
  \dcref{ch9:3.9}
  The theorem remains true under the weaker hypothesis that $f$ is a conformal
  diffeomorphism. It is not known whether the theorem is still true if $f$ is merely
  a diffeomorphism. This would imply that $S^2\times S^2$ does not carry a metric of
  positive curvature, since the map $f$ which is the antipodal map on each factor
  preserves the orientation (each factor reverses the orientation) and does not have
  a fixed point. For more details, see A. Weinstein [We 2].
\end{remark}

\begin{corollary}[Synge]
  \label{cor:dc-ch9-3-10}
  \dcref{ch9:3.10}
  \uses{thm:dc-ch9-3-7, cor:dc-ch9-3-3}
  Let $M^n$ be a compact manifold with positive sectional curvature.
  - a) If $M^n$ is orientable and $n$ is even, then $M$ is simply connected.
  - b) If $n$ is odd, then $M^n$ is orientable.
\end{corollary}
\begin{proof}
  \textbf{a)} Let $\pi:\tilde M\to M$ be the universal covering of $M$. Introduce on
  $\tilde M$ the covering metric, and orient $\tilde M$ so that $\pi$ preserves
  orientation. Because $M$ is compact with positive curvature, $K\ge\delta>0$; since
  $\pi$ is a local isometry, the same curvature condition holds on $\tilde M$. Since
  $\tilde M$ is complete, by \cref{cor:dc-ch9-3-3}, $\tilde M$ is compact. Let
  $k:\tilde M\to\tilde M$ be a covering transformation, that is, $\pi\circ k=\pi$
  (see Massey [Ma], p. 159). Then $k$ is an isometry of $\tilde M$ which preserves
  orientation. Because $n$ is even, we can use the theorem to conclude that $k$ has a
  fixed point. But a covering transformation with a fixed point is the identity. It
  follows that the group of covering transformations of $\tilde M$ (isomorphic to the
  fundamental group of $M$; see Massey [Ma], p. 163) reduces to the identity.
  Therefore $M$ is simply connected.

  \textbf{b)} Suppose that $M$ is not orientable, and consider the orientable double cover
  $\overline M$ of $M$ (see Exercise 12 of Chap. 0). Introduce on $\overline M$ the
  covering metric. Since $\overline M$ is the double cover of a compact manifold,
  $\overline M$ is compact. Let $k$ be a covering transformation of $\overline M$,
  $k\ne\mathrm{id}$. Because $M$ is not orientable, $k$ is an isometry which reverses
  the orientation of $\overline M$. Since $n$ is odd, we can apply \cref{thm:dc-ch9-3-7} which
  guarantees that $k$ has a fixed point. Therefore $k=\mathrm{id}$, which is a
  contradiction. $\square$
\end{proof}

\begin{remark}[Remark 3.11]
  \label{rem:dc-ch9-3-11}
  \dcref{ch9:3.11}
  The real projective space $P^2(\mathbf{R})$ of dimension two, which is compact,
  non-orientable and not simply connected, is an example which shows the necessity of
  orientability in a) and odd dimension in b). On the other hand, $P^3(\mathbf{R})$
  which is compact, orientable and not simply connected is an example which shows the
  necessity of even dimension in (a).
\end{remark}
